%=======================================================================
%  Phase 3: Obstruction, Null, Alignment, and the Lie Structure
%  Insert into au_fukaya_unified after
%  \section{Phase 3 After the Backbone: Slow-Manifold Relaxation}
%  Requires: amsmath, amssymb, amsthm, booktabs (already in preamble)
%  Fig.~\ref{fig:disc-strand} also needs: \usepackage{graphicx,tikz} with
%  \usetikzlibrary{arrows.meta,positioning,calc}
%=======================================================================

\section{Phase 3 Algebraic Structure: Obstruction, Null, Alignment}
\label{sec:phase3-algebra}

This section reports a measurement programme carried out entirely inside
phase 3 (steps $120$--$200$ of a $200$-step run on the standard corpus,
$P = 4{,}330{,}240$ parameters, $\mathrm{VOCAB}=1017$, $\mathrm{nnz}=1347$).
Its purpose is to determine which algebraic structure, if any, the
seven-node decomposition
\[
  \mathcal{N} = \{\mathrm{EMB},\ \mathrm{LN},\ \mathrm{FF},\
  W_Q,\ W_K,\ W_V,\ W_O\}
\]
carries. The programme is organised so that every positive claim is stated
against an explicit null, since the central methodological finding is that
several earlier claims in this document were reported without one.

\subsection{Trajectory geometry of the phase}

Phase 3 begins where the block gradient-profile drift settles
($0.130$ at steps $21$--$40$, falling to $0.038$ by step $100$ and flat
thereafter). Within the phase the step size is constant at
$\|d_t\|\approx 0.60$ and the discrete curvature is constant at
$\kappa \approx 0.70$--$0.74$, giving radius of curvature $R \approx 1.35$--$1.46$
against $\|\theta\|\approx 93$, i.e.
\begin{equation}
  R/\|\theta\| \;=\; 0.015 .
\end{equation}
The tangent autocorrelation decays with no negative lobe at any lag
($k=1{:}\,0.863$, $k=4{:}\,0.513$, $k=8{:}\,0.271$, $k=16{:}\,0.058$,
$k=32{:}\,0.002$, $k=64{:}\,0.002$), so the turning is diffusive, not
rotational. There is no coherent frame precession and no $32$-step period.

\begin{observation}[Chord--path dichotomy]
The useful chord over the phase is $\|\Delta\| = 19.6 = 0.22\|\theta\|$,
which is $\sim 14 R$. The trajectory is therefore not a stiff arc, yet the
straight chord from $\theta_{120}$ to $\theta_{200}$ is a monotone descent
path (recovery $40\%,\,73\%,\,93\%,\,100\%$ at chord fractions
$0.25,\,0.5,\,0.75,\,1.0$). The basin is chord-convex while the path is
diffusive. Consequently the base coordinate cannot be arc length; it must be
chord fraction toward the attractor.
\end{observation}

\subsection{The random-walk law and the persistence order parameter}

Let $S_n$ denote the parameter displacement accumulated over $n$ steps. For a
driftless random walk $\cos(S_n,S_N) = \sqrt{n/N}$ identically. Measured on
$120\to200$:

\begin{center}
\begin{tabular}{rcc}
\toprule
$n$ & $\cos(S_n, S_{80})$ & $\sqrt{n/80}$ \\
\midrule
5  & 0.2968 & 0.250 \\
10 & 0.3807 & 0.354 \\
20 & 0.4919 & 0.500 \\
30 & 0.6040 & 0.612 \\
40 & 0.7042 & 0.707 \\
60 & 0.8829 & 0.866 \\
\bottomrule
\end{tabular}
\end{center}

Agreement is within $0.02$ at every $n$. Phase 3 is statistically
indistinguishable from a driftless walk at the $80$-step scale. This
retrospectively explains every failed extrapolation experiment in this
document: there is no persistent direction to estimate, and the ceiling on
any local probe is set by the walk rather than by the estimator.

Define the \emph{persistence excess}
\begin{equation}
  E(W) \;=\; \cos\!\big(S_{W/2},\,S_W\big) \;-\; \tfrac{1}{\sqrt{2}},
  \label{eq:persistence}
\end{equation}
which vanishes identically for a driftless walk, giving a parameter-free null.

\begin{center}
\begin{tabular}{rccc}
\toprule
$W$ & $E$ on $40\to120$ & $E$ on $120\to200$ & s.e.m. \\
\midrule
4  & $+0.2523$ & $+0.2491$ & 0.002 \\
8  & $+0.2103$ & $+0.2098$ & 0.002 \\
16 & $+0.1510$ & $+0.1448$ & 0.003 \\
32 & $+0.1325$ & $+0.0739$ & 0.003 \\
64 & $+0.1621$ & $+0.0117$ & 0.003 \\
\bottomrule
\end{tabular}
\end{center}

The phases are indistinguishable at short windows (that part is optimizer
smoothing) and separate cleanly at $W=64$: $+0.162$ against $+0.012$, a
$\sim 50\sigma$ gap. $E(64)$ is therefore an order parameter for the
drift/walk transition, requiring three parameter snapshots and no gradient
evaluations. It replaces the hand-calibrated \texttt{drift\_thr} of the
compiler with a quantity whose null is analytic.

\begin{observation}[Gated extrapolation]
Applying $\theta \leftarrow \theta_t + \beta(\theta_t - \theta_{t-40})$
followed by $k$ corrector steps: at $t=80$ ($E(64)=+0.16$), $\beta=0.5$ with
$k=10$ reaches a value requiring $20$ ordinary steps, a net saving of $10$;
at $t=120$ the saving is $0$; at $t=160$ ($E(64)=+0.01$) every configuration
loses. The jump alone is uphill ($0.3274 \to 0.3485$) and finishes below
baseline only after correction, confirming the lift-then-select mechanism.
Optimal $\beta = 0.5$, not $1.0$; $\beta = 2$ is catastrophic throughout.
\end{observation}

\subsection{Cellular sheaf on the phase: \texorpdfstring{$H^0=0$}{H0=0}}

Partition $120\to200$ into four cells of $20$ steps and let the stalk over a
cell be the local displacement subspace. Effective ranks are
$3.26,\,3.25,\,3.63,\,3.36$ (CV $5\%$): \emph{no rank jump occurs anywhere in
the phase}, so there is no index event in the interior and the family
$\{A_b\}$ is trivialisable there. Principal angles between adjacent cells at
$k=6$ give exactly one surviving direction ($35$--$41^\circ$) with the
remaining five at $82$--$90^\circ$, overlap $0.10$--$0.12$.

\begin{proposition}[No global section]
The intersection of all four cell subspaces is trivial: the mean projector
$\tfrac14\sum_i P_i$ has top eigenvalue $0.43$--$0.51$ for every $k$ tested,
so $\dim H^0 = 0$ at any tolerance above $\cos > 0.5$. The net displacement
projects onto each individual cell at $0.48$--$0.66$ but lies in no
intersection: it is a quasi-section, present in every stalk and exact in none.
\end{proposition}

The reduced dynamics therefore exists locally in every cell and globally
nowhere, which is the precise sense in which phase 3 is simultaneously
measurable and unpredictable.

\subsection{Node interactions and the sign flip}

Over $120\to200$, per-parameter cancellation is $62.5\%$. Single-node oracle
jumps recover: EMB $88.4\%$, LN $21.9\%$, FF $31.8\%$, $W_Q$ $6.4\%$,
$W_K$ $-5.3\%$, $W_V$ $23.9\%$, $W_O$ $29.0\%$, summing to $196.2\%$ against
a joint recovery of $100\%$.

\begin{observation}[Additivity sign flip]
On $40\to120$ the skeleton is \emph{super}-additive (singles sum $26\%$,
joint $76\%$); on $120\to200$ the nodes are \emph{sub}-additive (singles sum
$196\%$, joint $100\%$). The sign of the additivity defect is a
threshold-free phase discriminator.
\end{observation}

Pairwise synergies $S_{ij} = \mathrm{rec}(i,j)-\mathrm{rec}(i)-\mathrm{rec}(j)$
reproduce the architecture: $W_V$--$W_O = -35.6$ (multiplicative composition
in the output path, so joint motion overshoots), $W_Q$--$W_K = +9.1$
(composition in the score, so each needs the other), and $W_K$ is a pure
connector --- solo $-5.3\%$ but positive with every partner, edge sum $+72.5$.
Its contribution lives on edges, not on its vertex.

The Gram matrix of functional changes $\cos(\delta f_i, \delta f_j)$ exhibits
four near-orthogonal blocks: $\{\mathrm{EMB},\mathrm{LN}\}$ (internal $0.331$),
$\{W_Q,W_K\}$ ($0.576$), $\{W_V,W_O\}$ ($0.612$), $\{\mathrm{FF}\}$, with
cross-block overlaps $0.02$--$0.16$ and $\mathrm{EMB}\leftrightarrow$attention
at $0.02$. Eigenvalues $1.854,\,1.437,\,1.343,\,0.902,\,0.653,\,0.424,\,0.387$:
the pairing is nondegenerate. Functional overlap predicts interaction sign,
$\mathrm{corr}(\cos(\delta f_i,\delta f_j),\,S_{ij}) = -0.37$.

\subsection{The null, and what it does to the closure claim}

Each node acts almost exactly linearly on its own
($\|\delta f(1.0\,d_g)\| / 2\|\delta f(0.5\,d_g)\| \in [0.916, 1.007]$), but the
nodes do not superpose: $\|\delta f(\text{full}) - \sum_g \delta f(g)\|$ is
$49.1\%$ of $\|\delta f(\text{full})\|$. All non-additivity is therefore
bilinear, defining a symmetric product
\begin{equation}
  m(u,v) \;=\; f(\theta+u+v) - f(\theta+u) - f(\theta+v) + f(\theta).
\end{equation}
Since $m$ maps $P\otimes P \to F$ with $F \neq P$, closure is only meaningful
after lowering the index with the Jacobian adjoint,
$\tilde m = J^{*}\!\circ m : P\otimes P \to P$.

\begin{center}
\begin{tabular}{lrrr}
\toprule
generator set & mean in-span & vs.\ null & \\
\midrule
node cuts of $\Delta$ (same segment)      & $0.695\%$ & $4{,}299\times$ \\
random directions per node, norm-matched  & $0.007\%$ & $42\times$ \\
random directions, second seed            & $0.006\%$ & $38\times$ \\
node cuts of an independent segment       & $1.575\%$ & $9{,}742\times$ \\
\bottomrule
\end{tabular}
\end{center}

The null for a random direction against a $7$-dimensional span in
$\mathbb{R}^{4{,}330{,}240}$ is $7/P = 1.62\times10^{-6}$.

\begin{observation}[Alignment without closure]
The product is $\sim\!10^3$--$10^4$ times more aligned with its own generators
than chance, and is \emph{more} aligned when the generators come from an
independent segment ($9{,}742\times$) than from the segment used to build it
($4{,}299\times$), ruling out inheritance. Random directions with identical
node support give only $\sim\!40\times$, so the alignment is a property of
on-trajectory directions and not of the block decomposition. Nevertheless the
product does not close: $99.3\%$ of its energy lies outside the generator span.
\end{observation}

The verdict is scale-independent. Writing $A(t)=m_t/t^2$, the raw product at
$t=1,\tfrac12,\tfrac14$, the two-point Richardson $H=2A(t/2)-A(t)$ and the
three-point Richardson $H=\tfrac13A(t)-2A(t/2)+\tfrac83A(t/4)$ give in-span
fractions $0.69\%,\,0.68\%,\,0.67\%,\,0.67\%,\,0.67\%$ respectively
($4{,}299\times$ to $4{,}130\times$ null), with pairwise scaling ratios
$3.29$--$3.90$ against the ideal $4$. Higher-order contamination therefore does
not account for the escape. The connected triple obstruction
satisfies $\|m_3\|/\|m_2\| = 0.339$ and $\|m_3\|/\|m_1\| = 0.184$ --- the
expansion converges geometrically --- but only $20.7\%$ of $m_3$ lies in
$\mathrm{span}(m_2)$. Effective ranks by order are
$2.88 \to 5.75 \to 8.24$, i.e.\ $\sim\!2.5$ new directions per order: linear
growth, far below the free bound $7^k$, but never closing.

\subsection{\texorpdfstring{$m_1$}{m1} from the flow, and why the
\texorpdfstring{$A_\infty$}{A-infinity} relations fail}

Take $\Phi(\theta) = \theta - \eta\nabla L(\theta)$ with a fixed batch, and set
\begin{align}
  m_0 &= \Phi(\theta)-\theta, &
  m_1(v) &= \Phi(\theta+v)-\Phi(\theta)-v, \\
  m_2(u,v) &= \Phi(\theta{+}u{+}v)-\Phi(\theta{+}u)-\Phi(\theta{+}v)+\Phi(\theta),
  \nonumber
\end{align}
with $m_3$ the connected third difference. \emph{Numerical warning:}
differencing $\Phi$ directly cancels $\|\theta\|\approx 93$ down to $10^{-5}$
and is pure single-precision roundoff --- diagnosed by $m_3$ scaling as $1.08$
instead of the required $8$. Differencing the gradient (where the $\theta$
terms cancel analytically) in double precision restores validity: $m_2$ scales
$3.32$ (expect $4$), $m_3$ scales $12.04$ (expect $8$).

With valid numerics both curved relations fail:
\begin{center}
\begin{tabular}{lrrr}
\toprule
relation & scale & residual & relative \\
\midrule
$n{=}1$: $m_1^2x + 2m_2(m_0,x)$, $x=\mathrm{EMB}$ & $4.64{\times}10^{-7}$ & $4.46{\times}10^{-7}$ & $0.961$ \\
$n{=}1$: $x=\mathrm{FF}$ & $1.69{\times}10^{-6}$ & $1.45{\times}10^{-6}$ & $0.856$ \\
$n{=}3$ Stasheff, EMB-FF-$W_Q$ & --- & --- & $\cos = 0.005$ \\
$n{=}3$ Stasheff, $W_Q$-$W_K$-$W_V$ & --- & --- & $\cos = -0.287$ \\
\bottomrule
\end{tabular}
\end{center}

\begin{proposition}[No $A_\infty$ structure in the gradient flow]
For $\Phi$ a gradient map, every $m_k$ is a derivative of the scalar potential
$L$: $m_1 = -\eta H$, $m_2 = -\eta\nabla^3 L$, $m_3 = -\eta\nabla^4 L$. These are
fully symmetric tensors. A fully symmetric family admits no associativity to
correct, so the associator is not the boundary of anything and the Stasheff
relations cannot hold. The measured failure is a consequence of the category,
not of the measurement.
\end{proposition}

This closes the symmetric route. Any antisymmetric structure must therefore
arise from an operation that is not a derivative of $L$.

\subsection{The Lie structure: node-restricted fields nearly close}

Let $P_i$ be the coordinate projector onto node $i$ and
\begin{equation}
  V_i(\theta) \;=\; -\,P_i\nabla L(\theta),
  \qquad
  [V_i,V_j] \;=\; -P_j H V_i + P_i H V_j .
  \label{eq:bracket}
\end{equation}
Although $H$ is symmetric, $P_iHP_j$ and $P_jHP_i$ occupy \emph{different}
blocks and cannot cancel; the bracket is nonvanishing. Antisymmetry is
structural and the Jacobi identity is automatic for vector fields, so the only
substantive question is closure. All $21$ brackets follow from $8$ gradient
evaluations.

\begin{center}
\begin{tabular}{lrrr}
\toprule
& symmetric product & Lie bracket \\
\midrule
mean in-span                & $0.7\%$ & $\mathbf{53.3\%}$ \\
proper null                 & $1.62{\times}10^{-6}$ & $2.62{\times}10^{-6}$ \\
ratio to null               & $4{,}299\times$ & $\mathbf{351{,}064\times}$ \\
effective-rank flow         & $3.76 \to 5.50$ (expands) & $2.40 \to 1.82$ (contracts) \\
energy outside span         & $99.3\%$ & $36.4\%$ \\
\bottomrule
\end{tabular}
\end{center}

Brackets are $0.469\times$ the magnitude of the generators, so this is a
leading-order structure. The control is decisive: projecting the \emph{same}
brackets onto random node-supported fields of matched norm collapses the
in-span fraction to $0.0078\%$ ($28.8\times$ null), a signal-to-control ratio
of $6{,}879$.

\begin{observation}[Contraction versus expansion]
The symmetric product expands rank and escapes almost completely; the
antisymmetric bracket contracts rank ($1.82 < 2.40$) and retains half its
energy in the generator span, escaping in only $\sim\!2.07$ effective
directions. This is the behaviour of a derived subalgebra
$[\mathfrak{g},\mathfrak{g}]\subseteq\mathfrak{g}$, missing closure by about
two directions rather than by everything.
\end{observation}

\subsection{Level two: the tower closes better than it opens}

The bracket fields $[V_i,V_j]$ are themselves vector fields, so the derived
series is computable. Let the level-1 span be
$\mathcal{G}_1=\mathrm{span}\{V_i\}\oplus\mathrm{span}\{P_j\,\mathrm{D}g_i\}$
($56$ dimensions, a superset of $\mathrm{span}\{V,[V,V]\}$), and evaluate
$[V_k,[V_i,V_j]]$ by finite differences of the bracket field.

\begin{center}
\begin{tabular}{lrrrr}
\toprule
$[V_k,[V_i,V_j]]$ & $\|\cdot\|$ & in $\mathrm{span}\{V\}$ & in $\mathcal{G}_1$ & vs.\ null \\
\midrule
$[\mathrm{EMB},[W_Q,W_K]]$ & $0.02814$ & $18.4\%$ & $91.3\%$ & $70{,}561\times$ \\
$[\mathrm{FF},[W_Q,W_K]]$  & $0.01445$ & $19.1\%$ & $91.5\%$ & $70{,}731\times$ \\
$[W_Q,[W_Q,W_K]]$          & $0.00140$ & $16.3\%$ & $90.7\%$ & $70{,}135\times$ \\
$[\mathrm{EMB},[W_V,W_O]]$ & $0.03162$ & $6.2\%$  & $39.5\%$ & $30{,}518\times$ \\
$[\mathrm{FF},[W_V,W_O]]$  & $0.01980$ & $10.5\%$ & $58.1\%$ & $44{,}918\times$ \\
$[W_Q,[W_V,W_O]]$          & $0.00121$ & $9.2\%$  & $50.2\%$ & $38{,}813\times$ \\
\midrule
mean                       &           & $12.8\%$ & $67.7\%$ & $52{,}342\times$ \\
\bottomrule
\end{tabular}
\end{center}

\begin{observation}[Root-space behaviour of $W_Q,W_K$]
The pair that fails hardest to close at level one closes hardest at level two.
$[W_Q,W_K]$ retains only $34.1\%$ in its own span, but every double bracket
$[V_k,[W_Q,W_K]]$ returns to $\mathcal{G}_1$ at $90.7$--$91.5\%$, and does so
uniformly across $k$. Conversely $[W_V,W_O]$, the most closed pair at level
one, generates level-two elements that escape ($39.5$--$58.1\%$). A generator
pair whose bracket leaves the span but whose adjoint orbit returns to it is the
behaviour of a root vector: $W_Q,W_K$ generate a new graded piece which is then
stable under the adjoint action, while $W_V,W_O$ behave as a solvable block that
does not.
\end{observation}

\subsection{The algebra is a stalk, not a global object}

Computing the brackets at step $120$ and projecting them onto the fields
$V_i(\theta_{c})$ at other checkpoints tests whether the closure is a property
of the phase or of the base point.

\begin{center}
\begin{tabular}{lrrrr}
\toprule
fields from & mean in-span & min & max & vs.\ null \\
\midrule
step $40$          & $0.11\%$  & $0.03$  & $0.24$  & $678\times$ \\
step $80$          & $0.14\%$  & $0.00$  & $0.36$  & $871\times$ \\
step $120$ (own)   & $44.97\%$ & $19.97$ & $58.08$ & $278{,}197\times$ \\
step $160$         & $1.53\%$  & $0.35$  & $3.90$  & $9{,}471\times$ \\
random node fields & $0.00\%$  & $0.00$  & $0.00$  & $3\times$ \\
\bottomrule
\end{tabular}
\end{center}

\begin{proposition}[Locality of the Lie structure]
Closure falls by a factor of $\sim\!30$ at $+40$ steps and $\sim\!300$ at
$-40$ steps, while remaining $10^2$--$10^4$ times above the random control
($3\times$). The Lie structure is therefore a local object with a decaying, not
vanishing, transport between checkpoints, and it is asymmetric in time: the
future ($9{,}471\times$) is better aligned than the past ($871\times$),
consistent with convergence toward the attractor.
\end{proposition}

This forces the algebraic reading into the sheaf-theoretic frame already
established in \S\ref{sec:phase3-algebra}: what exists is not one Lie algebra
but a \emph{field of Lie algebras over the base}, one per stalk, with weak
restriction maps. The measured transport ($0.034\times$ own over $40$ steps)
is quantitatively consistent with the principal-angle result that adjacent
cells share one direction out of six. The vanishing of $H^0$ for the transition
laws and the locality of the bracket closure are two readings of the same fact.

As a robustness note, the own-checkpoint closure evaluated at two independently
reached step-$120$ states gives $44.97\%$ and $53.31\%$, so the figure is
stable to about $\pm 8$ points across runs.

\subsection{Skeleton \texorpdfstring{$\to$}{->} dense: the transition, measured}

Evaluating \eqref{eq:bracket} at two checkpoints:

\begin{center}
\begin{tabular}{lrrrrr}
\toprule
subalgebra & $\|[\cdot,\cdot]\|_{80}$ & $\|[\cdot,\cdot]\|_{160}$ & growth
& in-span$_{80}$ & in-span$_{160}$ \\
\midrule
all seven                 & $0.00838$ & $0.03828$ & $4.57\times$ & $36.6\%$ & $55.9\%$ \\
skeleton EMB,LN,FF        & $0.02456$ & $0.13420$ & $5.47\times$ & $31.8\%$ & $63.0\%$ \\
attention $Q,K,V,O$       & $0.00152$ & $0.00490$ & $3.22\times$ & $31.8\%$ & $51.7\%$ \\
$W_Q,W_K$ only            & $0.00047$ & $0.00100$ & $2.11\times$ & $15.6\%$ & $34.1\%$ \\
$W_V,W_O$ only            & $0.00534$ & $0.01424$ & $2.67\times$ & $50.5\%$ & $60.5\%$ \\
\bottomrule
\end{tabular}
\end{center}

Every bracket grows and every subalgebra closes better as the run proceeds.
This is the skeleton-to-dense transition expressed algebraically: while the
skeleton carries the updates the Lie structure is weak and open; as attention
comes online the brackets strengthen by $2$--$5\times$ and closure rises from
$\sim\!32\%$ to $\sim\!63\%$.

\begin{observation}[$W_Q,W_K$ generate; $W_V,W_O$ close]
$W_V,W_O$ is the most closed pair at both checkpoints ($50.5\% \to 60.5\%$),
consistent with their multiplicative composition and functional overlap
$0.612$: they behave as a nearly solvable piece. $W_Q,W_K$ is the least closed
($15.6\% \to 34.1\%$), with
\[
  [V_Q,V_K] = +0.0203\,V_Q - 0.0140\,V_K + r,
  \qquad \|r\|/\|[V_Q,V_K]\| = 0.812
\]
at step $160$, and normalised bracket strength rising $1.263 \to 1.696$. A
generator pair whose bracket lands overwhelmingly outside its own span, with
opposite-sign coefficients, is the signature of a rank-2 system generating new
root spaces rather than of a closed subalgebra.
\end{observation}

\subsection{The layer grading and the culling mechanism}

Grading by depth rather than by parameter role sharpens every result. Let
$V_\ell = -P_\ell\nabla L$ with $\ell$ ranging over $\mathrm{EMB}$ and the six
blocks. Bracket closure rises to $62.2\%$ in-span ($384{,}812\times$ null),
above the $53.3\%$ obtained from the role decomposition, and decays
monotonically with depth separation:

\begin{center}
\begin{tabular}{rrrr}
\toprule
$|\ell-m|$ & pairs & mean in-span & mean $\|[\cdot,\cdot]\|$ \\
\midrule
1 & 5 & $63.7\%$ & $0.03117$ \\
2 & 4 & $62.1\%$ & $0.03024$ \\
3 & 3 & $59.8\%$ & $0.03099$ \\
4 & 2 & $57.7\%$ & $0.02952$ \\
5 & 1 & $56.8\%$ & $0.02373$ \\
\midrule
$\mathrm{EMB}$--$*$ & 6 & $64.6\%$ & $0.1127$ \\
\bottomrule
\end{tabular}
\end{center}

The mechanism is a rank collapse. Probing block $\ell$ with all seven
independent directions $\widehat V_m$ and collecting the responses
$P_\ell H\widehat V_m$:

\begin{center}
\begin{tabular}{lrrl}
\toprule
block & eff-rank & rank$(5\%)$ & normalised spectrum \\
\midrule
$\mathrm{EMB}$ & $1.07$ & $1$ & $1.00\ \ 0.02\ \ 0.01\ \ 0.00\ \ 0.00$ \\
$L_0$ & $1.10$ & $1$ & $1.00\ \ 0.03\ \ 0.01\ \ 0.01\ \ 0.00$ \\
$L_2$ & $1.11$ & $1$ & $1.00\ \ 0.02\ \ 0.02\ \ 0.01\ \ 0.01$ \\
$L_5$ & $1.13$ & $1$ & $1.00\ \ 0.04\ \ 0.01\ \ 0.00\ \ 0.00$ \\
\bottomrule
\end{tabular}
\end{center}

\begin{proposition}[Rank-one curvature response]\label{prop:rank-one}
Within every block the second-order response is rank one to within
$2$--$4\%$: $P_\ell H v \approx u_\ell\,(w_\ell\!\cdot\! v)$ for a fixed unit
$u_\ell$. This is the Jacobian structure of LayerNorm (a projection removing
mean and scale) composed with softmax (whose Jacobian $\mathrm{diag}(p)-pp^{\!\top}$
annihilates the constant direction).
\end{proposition}

The proposition explains the algebraic results rather than accompanying them.
Bracket closure is a single cosine per block, not a subspace overlap, which is
why it sits near $60\%$ rather than near $0$ or $1$; and adjoint stability at
level two is automatic, since once a bracket lies along $u_\ell$ every further
bracket does too. The $90.7$--$91.5\%$ closure of the $W_Q,W_K$ orbit is the
culling acting after the direction has been established.

\subsection{The obstruction to exploiting it}

Rank-one response invites a reduced Newton step: solve $g + Hd = 0$ on the
$K$ block coefficients, at a cost of $K{+}1$ gradient evaluations. Measured
against AdamW from the same state:

\begin{center}
\begin{tabular}{lrrrl}
\toprule
step & out-of-$u_\ell$ residual & $\mathrm{cond}(C)$ & equivalent AdamW steps & verdict \\
\midrule
$40$  & $72.0\%$ & $49.6$ & $1$ & no gain \\
$80$  & $61.9\%$ & $1.55{\times}10^{3}$ & $1$ & no gain \\
$120$ & $70.8\%$ & $477$ & $4$ & no gain \\
$160$ & $58.6\%$ & $1.47{\times}10^{3}$ & $1$ & no gain \\
\bottomrule
\end{tabular}
\end{center}

\begin{observation}[Fredholm obstruction, quantified]
Rank-one response means the image of the block-restricted curvature operator
is one-dimensional, so its cokernel is of codimension one in each block.
Between $59\%$ and $72\%$ of the gradient lies in that cokernel. By
Corollary~\ref{cor:fredholm} (the Fredholm alternative) the equation is
unsolvable for that component by any block-restricted update, and the reduced
Newton step correspondingly buys nothing. The culling that renders the algebra
tractable is the same structure that renders the gradient unreachable.
\end{observation}

What the rank-one result does justify is the phase-2 cull already validated
empirically: if $P_\ell H v \approx u_\ell(w_\ell\!\cdot\!v)$ to $2$--$4\%$, a
\emph{scalar} step size per block is not an approximation but the correct
curvature model, and maintaining a full second moment on those blocks is
wasted state. This is the mechanism behind $W_V,W_O \to$ scalar SGD and
$W_Q,W_K \to$ frozen diagonal being free within noise. It yields a
generalisable test: measure the response spectrum per block and cull those
exhibiting the $1.00 / 0.02$ gap. The depth decay above adds a second
prediction: inter-block coupling weakens with $|\ell-m|$, so layer-local
optimizer state and pipeline-parallel schedules are cheap exactly in
proportion to the measured off-diagonal of the coupling matrix $C$.

\subsection{Forward/backward alignment does not occur}

If the forward map (inter-token meaning) and the backward map (loss-driven
correction) came into alignment as training converged, the gradient would
progressively enter the reachable image and the out-of-$u_\ell$ residual would
decay to zero. Measured across a full run at $20$-step resolution:

\begin{center}
\begin{tabular}{rrrrrr}
\toprule
step & val & $\|g\|$ & mean residual & eff-rank & $\mathrm{cond}(C)$ \\
\midrule
$20$  & $2.1348$ & $0.5689$ & $76.7\%$ & $1.99$ & $44.3$ \\
$60$  & $0.6466$ & $0.3165$ & $67.5\%$ & $1.36$ & $87.8$ \\
$100$ & $0.1980$ & $0.1982$ & $61.9\%$ & $1.18$ & $189$ \\
$140$ & $0.0856$ & $0.2891$ & $69.4\%$ & $1.11$ & $3.44{\times}10^{4}$ \\
$160$ & $0.0704$ & $0.2010$ & $88.4\%$ & $1.25$ & $229$ \\
$200$ & $0.0768$ & $0.2221$ & $53.9\%$ & $1.08$ & $778$ \\
$220$ & $0.0985$ & $0.3717$ & $80.9\%$ & $1.09$ & $4.38{\times}10^{4}$ \\
\bottomrule
\end{tabular}
\end{center}

\begin{proposition}[The obstruction is permanent]
The mean residual over the first half of the run is $70.1\%$ and over the
second half is $70.1\%$; a linear fit gives slope $-0.0049\%$ per step,
extrapolating to zero at step $14{,}401$. The run attains its loss minimum at
step $160$ with the residual at its \emph{largest} value ($88.4\%$). Forward
image and backward demand therefore do not align, and the Fredholm obstruction
of the previous subsection is a permanent feature rather than a transient of
early training.
\end{proposition}

One quantity is monotone. The effective rank of the block response falls
$1.99 \to 1.64 \to 1.36 \to 1.11$ over steps $20$--$80$ and remains in
$[1.08, 1.25]$ thereafter. The rank-one culling is thus not a fixed property of
the architecture but \emph{crystallises} during skeleton formation, completing
at the skeleton-to-dense boundary. This dates the mechanism of
Proposition~\ref{prop:rank-one} and explains why the algebraic closure results
strengthen after step $80$.

\subsection{The projection path and its resolution scale}

The directions $u_\ell$ define rank-one projections $e_\ell = u_\ell u_\ell^{\!\top}$.
Two projections at principal angle $\vartheta$ satisfy
$\|e-f\| = \sin\vartheta$, and $\|e-f\| < 1$ is the standard criterion for
them to determine the same $K_0$ class. Sampling every step over a window in
phase 3:

\begin{center}
\begin{tabular}{rrrr}
\toprule
lag (steps) & mean $\cos$ & $\|e-f\|$ & \\
\midrule
$1$  & $0.891$ & $0.454$ \\
$2$  & $0.733$ & $0.680$ \\
$4$  & $0.402$ & $0.916$ \\
$8$  & $0.074$ & $0.997$ \\
\bottomrule
\end{tabular}
\end{center}

\begin{observation}[Continuity is a matter of sampling]
At $1$-step spacing $\|e-f\| = 0.454$, well inside the criterion, so the
projections form a continuous path and the $K_0$ class is defined along it.
The decorrelation time is $\approx 4$ steps; sampling at $20$-step intervals
(as in the preceding subsection) lies far beyond it and spuriously suggests a
discontinuous path. The second singular value satisfies $S_2/S_1 = 0.204$, so
$u_\ell$ is separated from its successor by a factor of five and the path is
not an artefact of a degenerate spectrum.
\end{observation}

The path is nevertheless diffusive rather than rotational. The per-step angle
is $\arccos(0.891) = 27.1^\circ$, accumulating $\approx 434^\circ$ over sixteen
steps, while the net displacement over the same interval is
$\arccos(0.158) = 80.9^\circ$: a path-length to net-displacement ratio of
$5.4$. There is therefore a well-defined path of projections but no coherent
rotation, and a training run --- which does not close a loop --- carries no
winding. Extracting a $K_1$ element would require constructing a loop
deliberately, for instance by a cyclic schedule returning to the same basin,
and measuring the holonomy of the transport around it. That construction is
the natural next step for the operator-theoretic programme and has not been
attempted here.

\section{The Tile Complex: One Space, Two Orientations}
\label{sec:tile-complex}

The preceding sections measured algebraic structure on a decomposition by
parameter role. This section replaces that decomposition with a refinement
hierarchy and finds that the resulting object is operator-algebraic rather than
Lie-theoretic, with a topology that is forced by the construction.

\subsection{Construction}

Each layer's parameter vector is partitioned, per group, into overlapping tiles
along its flat index. The resolution ladder is assigned by depth,
\[
  \text{forward}\quad L_0{:}16,\ L_1{:}32,\ L_2{:}64,\ L_3{:}256,\ L_4{:}512,\ L_5{:}1024,
\]
with the backward assignment mirrored ($L_5{:}16$ through $L_0{:}1024$).
Overlap is gated by the persistence excess of \S\ref{sec:phase3-algebra}: $10\%$
while $E(64) > 0.05$, $25\%$ below. On a $200$-step run the gate fired
unprompted at step $176$, where $E(64)$ crossed from $+0.060$ to $+0.043$.
Forward tiles record the weights after each step; backward tiles record the
applied update $\theta_t - \theta_{t-1}$, both reduced to mean and $|\max|$.
Groups $\mathrm{LN},\mathrm{FF},\mathrm{EMB}$ are tracked from step $1$ and
$W_Q,W_K,W_V,W_O$ from step $121$.

\begin{observation}[One complex, two orientations]
Forward and backward are not two spaces. They share the layer axis and the
index set and differ only in the orientation of the refinement maps: forward
branches toward the output, backward toward the input. The formal difference
$[\mathcal{H}_f]-[\mathcal{H}_b]$ is therefore taken on a single complex.
Ranks multiply to a constant at every depth,
\[
  16\cdot 1024 = 32\cdot 512 = 64\cdot 256 = 256\cdot 64
  = 512\cdot 32 = 1024\cdot 16 = 2^{14},
\]
while the graded difference $\operatorname{rank}\mathcal{H}_f -
\operatorname{rank}\mathcal{H}_b$ takes the values
$-1008,\,-480,\,-192,\,+192,\,+480,\,+1008$: antisymmetric about the middle,
summing to zero, and \emph{never} zero at any layer. The fixed point would be
$\sqrt{2^{14}} = 128$, which the ladder omits. The pairing is consequently
nondegenerate at every layer and the graded index changes sign between $L_2$
and $L_3$ by a jump rather than through a zero.
\end{observation}

Within a layer the tiles are overlapping intervals on a line, so the nerve of
the cover is a path graph, all intersections are contractible, and by the Nerve
Lemma each layer is homotopy equivalent to a point. All topology therefore
resides in the vertical direction. As a CW complex, per group and orientation:
$1904$ zero-cells, $1898$ horizontal and $1888$ vertical one-cells, and $875$
plaquettes, giving $\chi = -1007$.

\subsection{Stalks: a log power law in the loss}

Over each tile the stalk is the pair $(\alpha,\beta)$ fitted from
\begin{equation}
  \log\bigl|u_{\text{tile}}(t)\bigr|
  \;=\; \alpha\,\log\bigl(L(t)-H\bigr) \;+\; \beta,
  \qquad H = 0.062,
  \label{eq:powerlaw}
\end{equation}
where $u$ is the tile's $|\max|$ channel and $H$ is the corpus entropy floor.
Taking logarithms is what makes the stalk a vector space: addition and scalar
multiplication are defined, so restriction maps can be linear and cellular
sheaf cohomology is available. The metric on each stalk is supplied by the
diagonal Fisher information, computed with labels sampled from the model,
aggregated over the tile's index span.

\begin{center}
\begin{tabular}{lrrr|lrrr}
\toprule
\multicolumn{4}{c|}{FF forward (weights)} & \multicolumn{4}{c}{FF backward (updates)} \\
layer & tiles & $\alpha$ & $R^2$ & layer & tiles & $\alpha$ & $R^2$ \\
\midrule
$L_0$ & 16   & $-0.052$ & $0.729$ & $L_5$ & 16   & $+0.049$ & $0.733$ \\
$L_1$ & 32   & $-0.055$ & $0.776$ & $L_4$ & 32   & $+0.054$ & $0.697$ \\
$L_2$ & 64   & $-0.057$ & $0.765$ & $L_3$ & 64   & $+0.058$ & $0.658$ \\
$L_3$ & 256  & $-0.062$ & $0.774$ & $L_2$ & 256  & $+0.066$ & $0.596$ \\
$L_4$ & 512  & $-0.065$ & $0.782$ & $L_1$ & 512  & $+0.071$ & $0.556$ \\
$L_5$ & 1024 & $-0.068$ & $0.781$ & $L_0$ & 1024 & $+0.078$ & $0.530$ \\
\bottomrule
\end{tabular}
\end{center}

The exponents are mirror-symmetric between the two orientations, $\alpha_f
\approx -\alpha_b$, both monotone in depth: weights contract as the loss falls
while updates grow. $\mathrm{LN}$ and $W_K$ fit poorly ($R^2 = 0.02$--$0.32$),
so the power law in the loss is specific to $\mathrm{FF}$.

\subsection{Restriction maps as pants coproducts}

A refinement rung takes one parent tile to $r$ children with $r \in \{2,4\}$:
one input, several outputs. This is a coproduct
$\Delta : \mathcal{S}(p) \to \bigoplus_{i}\mathcal{S}(c_i)$, and the reverse
merge is the corresponding product. Fitting a single $2\times2$ map
$M : \mathcal{S}(p) \to \mathcal{S}(c)$ per rung by Fisher-weighted least
squares:

\begin{center}
\begin{tabular}{llrr}
\toprule
group/side & rung & $M$ & rel.\ residual \\
\midrule
FF bwd & $L_5\!\to\!L_4$ & $\left[\begin{smallmatrix}0.240 & -0.006\\ 0.598 & 1.008\end{smallmatrix}\right]$ & $0.001$ \\
FF bwd & $L_4\!\to\!L_3$ & $\left[\begin{smallmatrix}0.067 & -0.008\\ -0.251 & 1.003\end{smallmatrix}\right]$ & $0.001$ \\
FF bwd & $L_3\!\to\!L_2$ & $\left[\begin{smallmatrix}0.194 & -0.008\\ -0.470 & 1.004\end{smallmatrix}\right]$ & $0.002$ \\
FF bwd & $L_2\!\to\!L_1$ & $\left[\begin{smallmatrix}0.262 & -0.008\\ -1.300 & 0.992\end{smallmatrix}\right]$ & $0.003$ \\
FF bwd & $L_1\!\to\!L_0$ & $\left[\begin{smallmatrix}0.310 & -0.008\\ -1.503 & 0.990\end{smallmatrix}\right]$ & $0.003$ \\
\bottomrule
\end{tabular}
\end{center}

Residuals are $0.1$--$0.8\%$: the coproduct is essentially exact. Every map has
the same spectral shape, one eigenvalue near unity and one strongly contracted.
Across all rungs, groups and orientations the near-unit eigenvalue lies in
$[0.990,\,1.045]$ and the contracted one in $[0.005,\,0.53]$.

\subsection{The \texorpdfstring{$\beta$}{beta} line is a global section}

Comparing the eigenvalue-one eigenvector across every group and orientation:

\begin{center}
\begin{tabular}{lrrrrrr}
\toprule
 & FF/f & FF/b & LN/f & LN/b & $W_K$/f & $W_K$/b \\
\midrule
FF/f    & $1.0000$ & $0.9993$ & $0.9999$ & $0.9994$ & $0.9998$ & $0.9996$ \\
FF/b    & $0.9993$ & $1.0000$ & $0.9997$ & $1.0000$ & $0.9999$ & $1.0000$ \\
LN/f    & $0.9999$ & $0.9997$ & $1.0000$ & $0.9998$ & $1.0000$ & $0.9999$ \\
$W_K$/b & $0.9996$ & $1.0000$ & $0.9999$ & $1.0000$ & $1.0000$ & $1.0000$ \\
\bottomrule
\end{tabular}
\end{center}

\begin{proposition}[Universality of the fixed line within this cover]
All pairwise cosines exceed $0.9993$, and within-group agreement across rungs
exceeds $0.9985$. The eigenvalue-one eigenspace is a single line shared by
every stalk of every group in both orientations. Within the refinement system,
an observable is fixed by every restriction map exactly when it lies on this
line.
\end{proposition}

This is a statement about the tile refinement and not about the network:
\S\ref{sec:cover-dependence} shows that $\beta$ does not survive a change of
cover, so it should not be identified with a basic observable of the
optimisation dynamics in the sense of \S\ref{sec:phase3-algebra}.

\subsection{The inductive limit selects \texorpdfstring{$\beta$}{beta}}

Composing all five rungs:

\begin{center}
\begin{tabular}{lrrr}
\toprule
group/side & eigenvalues & $|\lambda_2|/|\lambda_1|$ & $\cos(\text{top}, \beta)$ \\
\midrule
FF/fwd    & $(0.0000,\ 1.1169)$ & $2\times10^{-5}$ & $1.0000$ \\
FF/bwd    & $(0.0002,\ 1.0257)$ & $2\times10^{-4}$ & $1.0000$ \\
LN/fwd    & $(0.0037,\ 1.2065)$ & $3\times10^{-3}$ & $1.0000$ \\
LN/bwd    & $(0.8615,\ 1.0909)$ & $0.790$          & $0.9998$ \\
$W_K$/fwd & $(0.0000,\ 1.0538)$ & $<10^{-5}$       & $1.0000$ \\
$W_K$/bwd & $(0.0001,\ 1.0435)$ & $1\times10^{-4}$ & $1.0000$ \\
\bottomrule
\end{tabular}
\end{center}

\begin{proposition}[$\beta$ is the Perron--Frobenius observable]
Five of six composites are rank one to within $10^{-5}$--$10^{-3}$, with
surviving eigenvector equal to the $\beta$ line at cosine $1.0000$. Refinement
discards every observable except $\beta$, which is therefore not merely basic
but the unique attracting observable of the refinement dynamics.
\end{proposition}

Two qualifications. The dominant eigenvalue exceeds unity everywhere
($1.026$--$1.207$), so the fixed \emph{line} is exact while the fixed
\emph{scale} is not: $\beta$ is weakly expanding under refinement. And
$\mathrm{LN}$ backward does not reduce, with
$|\lambda_2|/|\lambda_1| = 0.790$; it is also the group with the strongest
depth attenuation and the most tile-heterogeneous (CV $0.40$ against FF's
$0.095$). \S\ref{sec:reliability} shows this is not a fitting artefact and
identifies it as a second slow observable rather than an exception. Across two
seeds the eigenvalue \emph{ratio} replicates ($\mathrm{LN}$ $0.31$--$0.60$,
$\mathrm{FF}$ $0.0000$--$0.0155$) while the eigen\emph{vector} does not
($\mathrm{LN}_w$ gives $(+0.999,+0.051)$ then $(+0.491,+0.871)$), so the stable
object is the spectral gap and not the eigenspace. Since the stalk is
two-dimensional, the surviving second direction is $\alpha$ itself: for
$\mathrm{LN}$ the loss-response exponent survives refinement, for
$\mathrm{FF}$ it is annihilated. Determinant and $\lambda_2$ are not
independent evidence, as $\lambda_1 \approx 1.09$ throughout makes
$\det \approx 1.09\lambda_2$; no claim of an $SL_2$ or cluster-algebraic
structure is warranted.

\subsection{Validation: the collapse is not regression attenuation}
\label{sec:reliability}

The restriction maps of the previous subsections are least-squares fits of the
child stalk on the parent stalk. When a predictor is measured with error, the
fitted coefficient is biased toward zero by the reliability ratio
$\sigma^2_{\text{signal}}/(\sigma^2_{\text{signal}}+\sigma^2_{\text{noise}})$.
Since $\beta$ is a log-magnitude intercept (precisely estimated) while $\alpha$
is a slope fitted against $\log(L-H)$ (less so), attenuation alone would
produce a near-unit eigenvalue along $\beta$ and an eigenvalue equal to the
reliability of $\alpha$ along $\alpha$ --- reproducing the observed spectral
shape without any refinement geometry. This is a complete alternative
explanation of Propositions above and must be excluded.

Fitting $(\alpha,\beta)$ independently on odd and even steps and correlating
across tiles gives the split-half reliability, Spearman--Brown corrected.

\begin{center}
\begin{tabular}{lrrrr}
\toprule
group/side & $\mathrm{rel}(\alpha)$ & $\mathrm{rel}(\beta)$ & mean $\lambda_\alpha$ & $\lambda_\alpha - \mathrm{rel}(\alpha)$ \\
\midrule
FF/fwd    & $0.999$ & $1.000$ & $0.125$ & $-0.875$ \\
LN/fwd    & $0.999$ & $1.000$ & $0.392$ & $-0.607$ \\
$W_K$/fwd & $0.952$ & $0.983$ & $0.095$ & $-0.857$ \\
FF/bwd    & $0.887$ & $0.975$ & $0.208$ & $-0.680$ \\
LN/bwd    & $0.961$ & $0.994$ & $0.915$ & $-0.045$ \\
$W_K$/bwd & $0.584$ & $0.627$ & $0.318$ & $-0.266$ \\
\bottomrule
\end{tabular}
\end{center}

\begin{proposition}[The contraction is geometric, not statistical]
Over $30$ rungs, mean $\mathrm{rel}(\alpha) = 0.897$ against a mean contracted
eigenvalue of $0.342$; the correlation between reliability and eigenvalue is
$+0.20$ and the mean absolute gap is $0.571$. Under the attenuation null these
two quantities would coincide. In the forward bundles $\alpha$ is estimated at
reliability $0.95$--$0.999$ and is nonetheless contracted to $0.10$--$0.39$, so
refinement discards a well-measured coordinate. The collapse onto $\beta$ is a
property of the refinement system.
\end{proposition}

Two consequences for the preceding results.

First, $\mathrm{LN}$ backward is promoted from exception to finding. Its
$\alpha$ is measured at reliability $0.961$ --- comparable to the forward
bundles --- yet its restriction eigenvalue is $0.915$, and three of its five
rungs exceed unity ($1.016$, $1.045$, $1.017$). Refinement does not contract
$\alpha$ there; it mildly expands it. Under the attenuation null the least
reliable branch should contract most, so this runs opposite to the null in the
informative direction. $\mathrm{LN}$ backward therefore carries a second slow
observable that $\beta$ does not capture, and the inductive system is not
completely reducible on that branch.

Second, $W_K$ backward must be excluded rather than interpreted. Its split-half
reliabilities are $0.584$ and $0.627$, with one rung at $-0.02$ --- the two
halves are uncorrelated. Its stalks are unfitted parameters. This is the same
group whose Fisher state is worst conditioned ($\max/\min \approx 5\times10^{6}$,
normalised entropy $0.667$), and the two facts are plausibly one: Fisher mass
concentrated on a few tiles leaves the remainder with no signal to fit. The
$\beta$-line consistency table above accordingly rests on five reliable columns
rather than six, which does not affect its conclusion since those five agree at
cosine $\geq 0.9993$.

\subsection{The observable algebra is AF, and \texorpdfstring{$K_1$}{K1} vanishes by construction}

The branching ratios $2,2,4,2,2$ make the ladder a Bratteli diagram
$\mathbb{Z}^{16}\to\mathbb{Z}^{32}\to\mathbb{Z}^{64}\to\mathbb{Z}^{256}
\to\mathbb{Z}^{512}\to\mathbb{Z}^{1024}$, each edge of multiplicity one and
each child having a unique parent. The strand space is the path space of this
diagram: a Cantor set carrying the $2$-adic ultrametric
$d(s,s') = 2^{-\mathrm{lca}(s,s')}$ truncated at depth ten, with $2^{10}$
strands. The associated algebra is commutative AF, $C(X)$ with $X$ totally
disconnected.

\begin{proposition}[Forced vanishing of $K_1$]
For any AF algebra $K_1 = 0$. Every negative winding result reported in this
document is therefore a consequence of the observable algebra rather than of
the optimisation dynamics: the diffusive projection path, the absence of
coherent rotation, and the failure to extract a $K_1$ element were determined
before measurement. The claim is that \emph{the AF algebra generated by the
refinement tree has trivial $K_1$}, not that optimisation has trivial $K_1$;
enlarging the algebra, for instance by a crossed product with the optimisation
dynamics or by adjoining time as a circle, can restore it.
\end{proposition}

What survives is $K_0$, the dimension group of locally constant integer
functions, which is purely combinatorial, together with the state on it, which
is geometric. This is the $\alpha$/$\beta$ separation in operator-algebraic
form: $\alpha_f + \alpha_b \approx 0$ is virtual cancellation, visible to
$K$-theory and equal to zero there, while $\beta$ is metric data that
$K$-theory forgets.

\subsection{Faithfulness and the weighted grading}

\begin{center}
\begin{tabular}{lrrrr}
\toprule
group & tiles & zero-weight & $\max/\min$ & entropy/max \\
\midrule
FF    & $1904$ & $0$ & $9.13\times10^{2}$ & $0.851$ \\
LN    & $1904$ & $0$ & $1.37\times10^{3}$ & $0.833$ \\
$W_K$ & $1904$ & $0$ & $5.39\times10^{6}$ & $0.667$ \\
$W_V$ & $1904$ & $0$ & $1.48\times10^{2}$ & $0.890$ \\
$W_Q$ & $1904$ & $0$ & $2.32\times10^{6}$ & $0.791$ \\
\bottomrule
\end{tabular}
\end{center}

No tile carries zero Fisher weight in any group, so the state has full support
and the AF algebra admits no order ideal annihilated by it. The state is
nevertheless ill-conditioned on the score pair: $W_K$ and $W_Q$ concentrate
their Fisher mass with condition numbers above $10^{6}$ and normalised entropy
$0.67$--$0.79$, against $0.85$--$0.89$ for $\mathrm{FF},W_V,W_O$.

Backward magnitude increases monotonically with depth in every group, with
$L_5/L_0$ ratios $1.25$ (FF), $1.38$--$1.39$ (attention), $1.83$ (LN). The
backward bundle is therefore attenuated toward the input, and the exchange
$J$ between orientations satisfies $J_w^2 = w$ rather than $J^2 = 1$. Measured
on FF the profile is $w = (0.800, 0.838, 0.875, 0.936, 0.970, 1.000)$ across
$L_0$ to $L_5$: a mild smooth deformation, so the $\pm1$ eigenspaces remain
nearly orthogonal and the Hodge decomposition survives with respect to the
Fisher-weighted inner product.

\subsection{$\beta$ is a magnitude field, and magnitude carries no value}
\label{sec:magdir}

$\beta$ is the intercept of $\log|u_{\text{tile}}|$: a field of magnitudes over
tiles, with no directional content. Whether any scheme driven by such a field
can generate useful updates is therefore bounded by how much of an update's
value resides in its magnitude profile. Taking the true $20$-step displacement
$d$ at a checkpoint and decomposing it:

\begin{center}
\begin{tabular}{lrrr}
\toprule
arm & step $80$ & step $120$ & step $160$ \\
\midrule
true update                                   & $100.0\%$ & $100.0\%$ & $100.0\%$ \\
magnitude only (in-tile directions randomised) & $-21.3\%$ & $-38.0\%$ & $+17.9\%$ \\
direction only (magnitudes flattened per layer) & $+82.9\%$ & $+154.7\%$ & $+224.3\%$ \\
sign only (one global scale)                   & $+91.2\%$ & $+163.9\%$ & $+205.3\%$ \\
random vector, norm matched                    & $-20.9\%$ & $-29.9\%$ & $-3.4\%$ \\
\bottomrule
\end{tabular}
\end{center}

Entries are the fraction of the true progress retained.

\begin{proposition}[Ceiling on magnitude-driven updates]
Preserving the exact per-tile magnitude profile while randomising directions
within tiles performs indistinguishably from a random vector of matched norm at
every checkpoint. Any update scheme driven by a per-tile magnitude field,
$\beta$ included, is therefore bounded above by random.
\end{proposition}

Direction alone retains or exceeds the full progress, and sign alone suffices.
Since $\|\operatorname{sign}(d)\cdot\overline{|d|}\|_2 = \|d\|_1/\sqrt{P} \leq
\|d\|_2$ by Cauchy--Schwarz, the sign step has \emph{smaller} norm than the
optimiser's own step while making more progress at two of three checkpoints.
We record this as an observation warranting separate investigation rather than
a claim: it is a single $20$-step application at one scale, not a training run.

\subsection{$\beta$ does not survive a change of cover}
\label{sec:cover-dependence}

The preceding sections compute $\beta$ on one cover. Whether $\beta$ is
intrinsic to the parameters or an artefact of that cover is decided by
computing the per-parameter ground truth --- the cover by singletons --- and
asking how well each cover's regional $\beta$ predicts it. On a
$128\times256$ block of $\mathrm{FF}$ weights at $L_2$:

\begin{center}
\begin{tabular}{lrrr}
\toprule
cover & regions & $\mathrm{corr}(\beta)$ & $\mathrm{sd}(\beta_{\text{region}})$ \\
\midrule
rows (neurons)      & $128$ & $0.256$ & $0.056$ \\
columns (inputs)    & $256$ & $0.188$ & $0.039$ \\
contiguous tiles    & $64$  & $0.176$ & $0.037$ \\
contiguous tiles    & $512$ & $0.314$ & $0.069$ \\
random partition    & $128$ & $0.052$ & $0.011$ \\
random partition    & $512$ & $0.121$ & $0.023$ \\
\midrule
\multicolumn{2}{l}{singleton ground truth} & $1$ & $0.226$ \\
\bottomrule
\end{tabular}
\end{center}

\begin{proposition}[Cover dependence of $\beta$]
Structured covers correlate with the per-parameter $\beta$ at $0.18$--$0.31$
against $0.05$--$0.12$ for random partitions, so a structured cover resolves
roughly three times the random baseline and the signal is not nil. But the
best cover attains $\mathrm{sd} = 0.069$ against a singleton $\mathrm{sd}$ of
$0.226$: approximately $9\%$ of the variance. Covers agree with one another
only when one refines the other ($0.56$--$0.82$ among rows and contiguous
tiles, which coincide under row-major flattening); covers partitioning along
different axes share nothing. $\beta$ is therefore predominantly a property of
the cover rather than of the parameters.
\end{proposition}

Reported $R^2$ values are negative ($-2.3$ to $-2.6$) for every cover, but this
is a Jensen artefact rather than a failure of prediction: regional $\beta$ is
fitted to $\log(\overline{|u|})$ while the ground truth averages $\log|u|$, and
$\log \overline{x} \geq \overline{\log x}$, so regional $\beta$ is biased
upward systematically. Correlation is the appropriate statistic.

Taken with \S\ref{sec:magdir}, these results close the optimisation branch of
the programme by two independent routes: $\beta$ carries no directional
information, and it is not intrinsic to the parameters. What survives is
narrower and should be kept distinct from $\beta$: the spectral gap of the
refinement operator ($\mathrm{LN}$ $0.31$--$0.60$ against $\mathrm{FF}$
$0.0000$--$0.0155$ across two seeds), the persistence excess $E(64)$, which is
computed from whole-parameter displacements with no cover at all, and the
rank-one curvature response, which is a property of the block projectors. None
of these depends on $\beta$ being intrinsic.

\subsection{Strands and the \texorpdfstring{$p$}{p}-adic metric}
\label{sec:strands}

A strand is a root-to-leaf path through the refinement ladder. With branchings
$2,2,4,2,2$ there are $2^{10}=1024$ strands, each identified by its leaf, and
each carrying the concatenated stalks $(\alpha,\beta)$ over six layers ---
a point in $\mathbb{R}^{12}$, $z$-scored per layer. The leaf set carries the
ultrametric $d(s,s') = p^{-\mathrm{lca}(s,s')}$.

Two remarks are needed before the measurements. First, the forward and backward
strand sets are not canonically paired: forward paths refine toward the output
while backward paths refine toward the input, so a forward strand and a
backward strand coincide at a layer only through the indexing convention. The
$\mathcal{H}_f/\mathcal{H}_b$ pairing must therefore be taken layer-wise, as it
is throughout. Second, $p^{-\mathrm{lca}}$ is a monotone function of the integer
$\mathrm{lca}$, so $p=2$, $3$ and $5$ induce identical rank orderings on the
true tree; the choice of $p$ is a labelling convention with no empirical
content. Measured Spearman correlations for $p = 2,3,5$ agree to three decimals.

\begin{center}
\begin{tabular}{lrrrr}
\toprule
metric (Spearman vs.\ strand distance) & FF/fwd & FF/bwd & LN/fwd & LN/bwd \\
\midrule
true tree ($p=2,3,5$ identical)  & $0.243$ & $0.243$ & $0.270$ & $0.307$ \\
$2$-adic on index                & $0.109$ & $0.178$ & $0.195$ & $0.193$ \\
$3$-adic on index                & $0.145$ & $0.239$ & $0.355$ & $0.249$ \\
$5$-adic on index                & $0.136$ & $0.269$ & $0.175$ & $0.137$ \\
$|i-j|$ (plain contiguity)       & $0.189$ & $0.237$ & $0.195$ & $0.092$ \\
permuted tree (shape kept)       & $0.001$ & $-0.003$ & $0.000$ & $-0.006$ \\
partial, tree given $\log|i-j|$  & $0.074$ & $0.061$ & $0.099$ & $0.196$ \\
\bottomrule
\end{tabular}
\end{center}

Strand distance does fall monotonically with shared-ancestor depth, with
shallow-to-deep ratios of $1.92$--$4.86$, so strands are not independent noise
vectors.

\begin{proposition}[The strand hierarchy is contiguity]
Retaining the tree's combinatorial shape while permuting which parameter index
occupies which leaf reduces the correlation from $0.24$--$0.31$ to
$0.000 \pm 0.006$ in all four branches. The hierarchy therefore contributes
nothing independently of its alignment with contiguous parameter blocks.
Arbitrary $p$-adic trees on the index perform comparably to the true refinement
and in one case better ($3$-adic, LN/fwd, $0.355$ against $0.270$). The
ultrametric structure of the strand space is inherited from the contiguity of
the tiling, not discovered in the network.
\end{proposition}

\subsection{What distinguishes LayerNorm, and what does not}
\label{sec:ln-status}

Four measurements appeared to single out $\mathrm{LN}$ under refinement: the
spectral gap, the near-unit determinant, the strand-distance ratio, and the
partial correlation of \S\ref{sec:strands}. They do not survive equally.

The determinant is not independent evidence: $\lambda_1 \approx 1.09$
throughout, so $\det \approx 1.09\lambda_2$ and the two are one measurement.

The strand residual fails decomposition. Combined $\mathrm{LN}$ gives partial
correlation $0.192$ and $0.208$ on two seeds, but splitting it yields
$\mathrm{LN}_w$ at $0.127/0.083$ and $\mathrm{LN}_b$ at $0.054/0.063$ --- both
below the whole. A quantity exceeding both its parts locates the structure in
the joining, and the combined group is
$\mathtt{attn.ln.weight}\,|\,\mathtt{attn.ln.bias}\,|\,\mathtt{ff.n.weight}\,|\,\mathtt{ff.n.bias}$:
four contiguous blocks alternating between gains near $1$ and offsets near $0$.
The residual is the concatenation order. Moreover $W_V$ reaches $0.103/0.135$,
matching or exceeding $\mathrm{LN}_w$, so no LayerNorm signature remains after
the split.

The spectral gap survives every control:

\begin{center}
\begin{tabular}{lrrrr}
\toprule
group & elements/tile & $\lambda_2/\lambda_1$ (seed 17) & (seed 23) \\
\midrule
$\mathrm{LN}_w$ & $0.5$ & $0.604$ & $0.516$ \\
$\mathrm{LN}_b$ & $0.5$ & $0.526$ & $0.330$ \\
$\mathrm{LN}$   & $1.0$ & $0.464$ & $0.313$ \\
$\mathrm{FF}$   & $384.5$ & $0.0008$ & $0.0155$ \\
$\mathrm{FF}1024$ & $1.0$ & $0.0000$ & $0.0001$ \\
$W_V$           & $64.0$ & $0.0000$ & $0.0000$ \\
\bottomrule
\end{tabular}
\end{center}

$\mathrm{FF}1024$, a $1024$-element slice of $\mathrm{FF}$ tiled with the
identical ladder, matches $\mathrm{LN}$'s occupancy exactly and behaves like
$\mathrm{FF}$, excluding occupancy. Both $\mathrm{LN}_w$ and $\mathrm{LN}_b$
show the gap independently, excluding concatenation. The separation is two to
four orders of magnitude and replicates across seeds.

\begin{observation}[The one surviving LayerNorm result]
Under this refinement, LayerNorm's transfer operator retains a second
eigenvalue of order $0.5$ while every other group annihilates it to order
$10^{-3}$. Since the stalk is two-dimensional, the surviving direction is
$\alpha$: the loss-response exponent survives refinement for $\mathrm{LN}$ and
is annihilated for $\mathrm{FF}$. This is consistent with LayerNorm's Jacobian
being a projection, $I - P_{\mathbf 1} - P_x$, and hence less dissipative than
an unconstrained linear map, but the mechanism is inferred rather than measured.
\end{observation}

\subsection{Summary, and what the construction does not establish}

The cover is the computational architecture; the latent foliation is prior to
it and is measured only through the cover, so cohomological statements
constrain the cover rather than the space. Stalks carry
$(\mathcal{H}_f,\mathcal{H}_b,D_{\text{sheaf}},\beta)$ with
$\mathcal{H}_f = \ell^2(\text{tiles}_{\text{fwd}})\otimes\mathbb{R}^2$ of
dimension $3808$ per group, Fisher supplying the inner product, and
$D_{\text{sheaf}} = \delta+\delta^{*}$ the Hodge--Dirac operator of the
restriction system --- distinct from the stochastic generator
$-G\cdot\nabla+\tfrac12\Sigma\!:\!\nabla^2$, which acts on functions over
parameter space and not on stalks.

Four of the five structural claims are theorems about the construction rather
than measurements of training. The hierarchy exists because it was built; the
maps form an inductive system because any nested partition does; $K_1=0$
follows from the Bratteli diagram having unique parents and finite branching;
and faithfulness fails only if the architecture carries structurally dead
parameters. The fifth, the Perron--Frobenius collapse, passes its null test
(\S\ref{sec:reliability}) and is a genuine property of the refinement system,
but $\beta$ then fails cover-independence (\S\ref{sec:cover-dependence}) and
carries no directional content (\S\ref{sec:magdir}). The strand ultrametric
reduces to contiguity (\S\ref{sec:strands}).

What survives controls is therefore narrow and should be stated without the
operator-algebraic framing, on which none of it depends: the persistence excess
$E(64)$, computed from whole-parameter displacements with no cover at all; the
rank-one curvature response, a property of the block projectors, which derives
the scalar-per-block preconditioner from measurement rather than convention;
the Fredholm obstruction, which explains why block-local methods cannot close
the gap; and the LayerNorm spectral gap of \S\ref{sec:ln-status}, which is a
fact about the refinement operator and not about the network.

The vanishing of the topological invariants is correspondingly not evidence
about optimisation. $K_1$ vanishes because an AF algebra admits no $K_1$; the
index is forced by dimension counting in finite dimensions; the strand
ultrametric is the tiling's own adjacency. Establishing whether any of this
describes training rather than the refinement requires a cover that is not a
nested partition, where the algebra need not be AF, and that has not been
attempted.

\section{The Sign--Magnitude Decomposition of the Update}
\label{sec:sign-magnitude}

The tile programme sought a reduced observable from which updates could be
generated. \S\ref{sec:magdir} and \S\ref{sec:cover-dependence} showed that
$\beta$ is not that observable. This section reports what the same experimental
apparatus finds when the question is inverted: rather than asking which reduced
quantity predicts the update, we ask how far the update can be compressed
before training degrades.

\subsection{Necessary motion, in hindsight}

Over $200$ steps the per-parameter path length exceeds the net displacement by
a factor of $2.3$: writing $\mathrm{net}_i = |\sum_t d_t^{(i)}|$ and
$\mathrm{path}_i = \sum_t |d_t^{(i)}|$, the aggregate ratio is
$\sum_i \mathrm{net}_i / \sum_i \mathrm{path}_i = 0.441$, so $55.9\%$ of all
motion cancels. Truncating the net displacement to its largest coordinates and
applying the result to $\theta_0$ gives:

\begin{center}
\begin{tabular}{lrrrr}
\toprule
top $k\%$ of $|{\rm net}|$ & $\%$ of $\sum|{\rm net}|$ & val & improvement kept & random $k\%$ control \\
\midrule
$1\%$   & $4.1\%$  & $4.9560$ & $-11.2\%$ & $0.7\%$ \\
$5\%$   & $16.0\%$ & $6.2832$ & $-41.5\%$ & $3.0\%$ \\
$10\%$  & $27.8\%$ & $5.6599$ & $-27.3\%$ & $5.6\%$ \\
$25\%$  & $53.9\%$ & $1.6289$ & $64.7\%$  & $5.3\%$ \\
$50\%$  & $81.4\%$ & $0.0920$ & $99.8\%$  & $5.5\%$ \\
$100\%$ & $100\%$  & $0.0754$ & $100.0\%$ & $99.7\%$ \\
\bottomrule
\end{tabular}
\end{center}

Roughly half the coordinates are required, and partial application is
\emph{worse than doing nothing} below $25\%$: the $5\%$ arm lands at val
$6.28$ against an initial $4.47$. Moving a LayerNorm gain without its paired
bias, or $W_O$ without $W_V$, leaves the network in a state that is neither the
initialisation nor the trained model --- the non-monotonicity is the signature
of the block coupling measured in \S\ref{sec:phase3-algebra} ($S_{W_V W_O} =
-35.6$). Combining the two factors, $0.441 \times 0.814 = 35.9\%$ of the total
motion is necessary in hindsight.

This necessary mass is not localised. Density (mass per parameter, normalised)
ranges only over $0.87$--$1.28$ across the seven components, with the embedding
highest; across the $1024$ strands the Gini coefficient is $0.095$ with lag-one
autocorrelation $+0.995$, so the mass field is a smooth plateau rather than a
concentrated chunk, and the top $10\%$ of strands carry $13.2\%$ of it. Forward
and backward ladders agree to three decimals. Nor does the corpus predict it:
per-token displacement in the embedding correlates with token frequency at
$-0.103$, with log frequency at $-0.142$, and with the number of distinct
successors at $-0.112$, varying by $5\%$ across frequency deciles.

\begin{observation}[Cancellation is not removable]
Rerunning from the same initialisation while zeroing, at every step, any update
component opposing $\operatorname{sign}(D_i)$ discards $41.4\%$ of the motion
and reaches val $0.0676$ against a baseline $0.0677$. But retaining
\emph{only} the opposing components --- $42.4\%$ of the motion --- also
converges, to val $0.0937$ ($99.4\%$). Any cull that leaves the optimiser free
to recompute gradients is not a constraint; the trajectory re-routes. The
$55.9\%$ cancellation is an accounting artefact of hindsight, not spendable
waste.
\end{observation}

\subsection{One constant per strand per step}

Replace the applied update $d$ at every step by a quantised version in which
each of $T$ contiguous chunks of the parameter vector receives a single
constant. Two schemes: retain the per-parameter sign and supply one magnitude
per chunk, $d' = \operatorname{sign}(d)\odot c_{\text{chunk}}$ with
$c = \overline{|d|}$; or make the chunk uniform, $d' = \bar d_{\text{chunk}}$.
Gradients are recomputed at each new point, so the trajectory is free to
diverge.

\begin{center}
\begin{tabular}{rrlrr}
\toprule
$T$ & numbers/step & scheme & val & improvement kept \\
\midrule
--- & $4{,}330{,}240$ & baseline (full update) & $0.0677$ & $100.0\%$ \\
$1$    & $1$    & $\operatorname{sign}(d)\times\overline{|d|}$ & $0.0646$ & $100.1\%$ \\
$64$   & $64$   & $\operatorname{sign}(d)\times\overline{|d|}$ & $0.0571$ & $100.2\%$ \\
$1024$ & $1024$ & $\operatorname{sign}(d)\times\overline{|d|}$ & $0.0572$ & $100.2\%$ \\
$1024$ & $1024$ & uniform signed constant & $4.4228$ & $0.9\%$ \\
\bottomrule
\end{tabular}
\end{center}

\begin{proposition}[Sign carries the update; magnitude is one number]
A single global scalar per step, applied with the per-parameter sign pattern,
retains $100.1\%$ of the loss improvement over a full $200$-step run --- a
$4.33\times10^{6}$ compression of the magnitude field with no degradation.
Making the constant uniform within a chunk, thereby discarding the
per-parameter sign, collapses training to $0.9\%$. The update therefore
decomposes into a $4.33$\,Mbit sign pattern carrying essentially all the
content and a magnitude field compressible to a single float.
\end{proposition}

The per-step payload falls from $139$\,Mbit to $4.33$\,Mbit plus $32$ bits, a
$32\times$ information compression that is exact rather than approximate. This
corroborates \S\ref{sec:magdir} from the opposite direction: there, retaining
magnitudes while randomising directions performed indistinguishably from a
random vector; here, retaining directions while flattening magnitudes to one
global constant loses nothing. Both experiments locate the entire content of an
update in its sign pattern, and the present one does so over a full training
run rather than a single application.

Two qualifications. The compression is of \emph{bandwidth}, not of computation:
the backward pass is still required to obtain the signs, and costs the same
whether one bit or thirty-two are retained per coordinate. The relevance is
therefore to distributed training, where gradient exchange is the bottleneck
and a one-bit-plus-scale payload is exactly what is transmitted; this is the
known $1$-bit SGD and signSGD family, and what is new here is only the
measurement that on this problem the scheme is not a lossy approximation but
marginally superior to AdamW ($0.0572$ against $0.0677$). Second, this is one
seed, one architecture, and a corpus consisting of a single repeated sentence,
a regime in which sign-based methods are known to be favoured. The exportable
claim is the decomposition, which two independent experiments establish, not
the superiority.

\subsection{Scope and limitations}

All measurements are at a single corpus (one $1364$-token base sequence,
$300$ loops), a single architecture, and one to four checkpoints per claim.
The tile complex of \S\ref{sec:tile-complex} rests on one $200$-step run; the
power-law stalks are fitted on the $|\max|$ channel only, and the restriction
maps are single $2\times2$ fits per rung rather than per parent, so their
$0.1$--$0.8\%$ residuals measure consistency of the rung, not of individual
tiles. Per \S\ref{sec:reliability}, the $W_K$ backward branch fails split-half
reliability and is excluded from interpretation throughout.

Four of the five structural claims about the tile complex are theorems about
the construction rather than measurements of training: the hierarchy exists
because it was built, the maps form an inductive system because any nested
partition does, $K_1=0$ follows from the Bratteli diagram having unique parents
and finite branching, and faithfulness fails only if the architecture carries
structurally dead parameters. Only the Perron--Frobenius collapse carries
empirical weight, which is why \S\ref{sec:reliability} tests it against its
null. Relatedly, $\beta$ and $\operatorname{tr}\Sigma$ should not be conflated:
the former is a log-magnitude intercept of tile updates, the latter a trace of
the gradient covariance, and the evidence for each is independent. Following
\S\ref{sec:magdir} and \S\ref{sec:cover-dependence}, $\beta$ should be
described as the dominant eigendirection of a particular refinement system --- a
reproducible fact about that system --- and not as an intrinsic observable of
the network. The cover-dependence test covers one block of one layer of one
group; extending it to $\mathrm{LN}$, where the refinement operator behaves
differently, is outstanding.

On the question the programme originally posed, the evidence assembled here is
negative and the negatives are independent. Phase 3 obeys the driftless-walk
law, so no local estimator can recover a persistent direction; the
block-restricted curvature operator leaves $59$--$72\%$ of the gradient in its
cokernel at every checkpoint, largest at the loss minimum, so no block-local
method can supply what is missing; and the reduced observable $\beta$ carries
magnitude only, which \S\ref{sec:magdir} shows is worth no more than a random
vector. Gradient descent computes direction and magnitude at the same cost, and
direction is the component that carries the value. The single measured lever is
the $E(64)$-gated jump, saving $+2$, $+14$ and $-2$ steps across three seeds:
positive in mean, unreliable per run, and the correct order of magnitude for a
regime governed by a random walk. The observation developed in \S\ref{sec:sign-magnitude} --- that the entire
content of an update lies in its sign pattern, with the magnitude field
compressible to a single scalar at no cost over a full run --- did not arise
from this framework and is the lead most worth pursuing.
The seven fields are masks of a single gradient vector and have effective rank
$2.40$ of $7$, so their mutual independence is modest. Finite differences use
$h=10^{-4}$ in double precision throughout; single precision is \emph{not}
sufficient and silently returns roundoff, as diagnosed above. The gated-jump
saving of $10$ steps is a single measurement on one seed, and validation noise
near the entropy floor is $\pm 20\%$; seed replication remains outstanding.
The level-two computation covers $9$ of the $147$ possible double brackets,
chosen to contrast the $W_Q,W_K$ and $W_V,W_O$ blocks, and a full enumeration
has not been performed.

The claims we regard as established are: (i) the phase-3 trajectory obeys the
driftless-walk law $\cos(S_n,S_N)=\sqrt{n/N}$ to within $0.02$, so no local
estimator can recover the chord, and $E(64)$ separates the phases at
$\sim\!50\sigma$; (ii) the symmetric product built from node displacements is
$\sim\!4{,}200\times$ better aligned with its generators than chance yet closes
at only $0.7\%$, a verdict stable under three-point Richardson extrapolation
and reproduced with generators from an independent segment; and (iii) the
node-restricted gradient fields carry a Lie structure closing at $45$--$53\%$
in-span, $\sim\!3\times10^{5}$ times its null and $\sim\!7\times10^{3}$ times
its matched control, which strengthens by $2$--$5\times$ and closes from
$\sim\!32\%$ to $\sim\!63\%$ across the skeleton-to-dense transition.

Three claims are refuted. The hypothesis that forward and backward maps come
into alignment as training converges is contradicted: the out-of-image residual
is flat at $70.1\%$ across both halves of a run and is largest at the loss
minimum. The reduced Newton step built on the rank-one response
gives no speedup at any checkpoint, because $59$--$72\%$ of the gradient lies
in the cokernel of the block-restricted curvature operator; the structure is
explanatory, not exploitable, as an optimizer. And the Lie structure is not a
single global algebra: cross-checkpoint closure decays by $\sim\!30\times$ over $40$ steps,
so the object is a sheaf of algebras over the base rather than one algebra,
and any Kac--Moody presentation would have to be stalk-local.

Two computations would decide the remaining structure. First, a full
enumeration of level-two brackets, testing whether the $90\%$ closure observed
for the $W_Q,W_K$ orbit holds generally or is specific to that pair; if it
holds, the graded pieces are stable under the adjoint action and structure
constants can be extracted stalk by stalk. Second, transport of the bracket
basis along the flow, to determine whether the $0.034\times$ decay over $40$
steps is genuine loss of structure or merely a rotating frame --- the
distinction between a sheaf with weak restriction maps and one whose stalks are
isomorphic under a connection that has not yet been identified.

\section{The Sign Process and the Stability Field}
\label{sec:sign-process}

The magnitude-versus-direction result of \S\ref{sec:sign-magnitude} left the
sign pattern as the carrier of the update. This section characterises that
pattern as a dynamical object and identifies the variable that organises it.
It is a record of experiments in the order run, including those that failed,
because the failures are what narrowed the object.

\subsection{Sign persistence and stale-sign at matched compute}
\label{sec:stalesign}

The step-to-step sign agreement of the applied update is $87.7\%$, flat across
the run ($88.6\%$ first fifty steps, $87.5\%$ last fifty, never below
$77.8\%$). This suggests reusing a cached sign to skip backward passes. Refreshing
the true sign every $k$ steps and applying $\operatorname{sign}\times$ scalar in
between:

\begin{center}
\begin{tabular}{lrrr}
\toprule
scheme & backward passes & val & vs.\ equal-compute baseline \\
\midrule
stale-sign $k=2$ & $100$ & $0.1670$ & $13\%$ better than $100$ full steps ($0.1927$) \\
stale-sign $k=4$ & $50$  & $0.6606$ & $25\%$ better than $50$ full steps ($0.8761$) \\
\bottomrule
\end{tabular}
\end{center}

\begin{observation}[The first compute-matched saving]
At equal backward passes, interleaving free stale-sign updates beats honest
AdamW by $13\%$ at $k=2$ and $25\%$ at $k=4$. This is the only lever in the
programme that survives a compute-matched control; the earlier gated jump and
all magnitude schemes do not. It works because it reuses a measured direction
rather than predicting one, sidestepping the driftless-walk obstruction.
\end{observation}

The sparse-event refinement of this idea fails. Sign flips are not rare:
$12.3\%$ of coordinates flip per step ($530{,}540$ of $4.33$M), flat across the
run, so ``detect only boundary crossings'' has no sparsity to exploit --- the
orthant label changes almost completely each step.

\subsection{Momentum is the transport projection}
\label{sec:momentum}

The oscillatory component of the update is separable and its separator already
exists: Adam's $\beta_1$ is a first-order low-pass filter. Sweeping it:

\begin{center}
\begin{tabular}{lrrrrr}
\toprule
$\beta_1$ & val & cancellation & path & $\cos(d_t,d_{t+1})$ & flip rate \\
\midrule
$0.00$ & $0.1122$ & $91.3\%$ & $599.1$ & $-0.093$ & $51.1\%$ \\
$0.90$ & $0.0677$ & $67.4\%$ & $160.5$ & $\phantom{-}0.906$ & $12.3\%$ \\
$0.99$ & $0.0682$ & $41.2\%$ & $116.9$ & $\phantom{-}0.987$ & $\phantom{0}4.5\%$ \\
\bottomrule
\end{tabular}
\end{center}

Every quantity moves monotonically with $\beta_1$. At $\beta_1=0$ the dynamics
is a genuine walk (consecutive steps \emph{anticorrelated}, half of all
coordinates flipping per step); the $\sqrt{n/N}$ law of
\S\ref{sec:phase3-algebra} is the filtered image of this. At $\beta_1=0.99$ the
path shortens $27\%$, sign flips fall to a quarter, and val is unchanged within
floor noise. Two consequences: the oscillation is real and largely removable
but removing it buys only a shorter path to the same endpoint; and the
$55.9\%$ cancellation figure quoted throughout is a property of Adam at
$\beta_1=0.9$, ranging $41$--$91\%$ with the filter, not a property of the loss
surface.

\subsection{Stochasticity is the regulariser, not coverage}
\label{sec:deterministic}

Suppressing the transverse motion by making the gradient deterministic
converges to the wrong place, and the reason is not batch coverage.

\begin{center}
\begin{tabular}{lrrrr}
\toprule
arm & val & $\cos(d,\text{to-final})$ & $\cos(d_t,d_{t+1})$ & path \\
\midrule
stochastic (baseline) & $0.0677$ & $0.356$ & $0.906$ & $160.5$ \\
fixed single batch    & $6.8069$ & $0.843$ & $0.996$ & $\phantom{0}68.2$ \\
deterministic full coverage & $5.9505$ & $0.847$ & $0.997$ & $\phantom{0}86.8$ \\
\bottomrule
\end{tabular}
\end{center}

Full coverage ($1344$ of $1364$ base tokens) reproduces the fixed-batch
geometry exactly and fails just as badly. Any deterministic gradient produces a
smooth $2.4\times$-shorter path that memorises: train loss $10^{-4}$ against val
$5.93$. The regimes separate at step $\sim\!20$. A rescue confirms the damage is
reversible: $50$ deterministic steps (landing at val $3.30$) followed by $150$
stochastic reach $0.0660$, \emph{beating} pure stochastic ($0.0794$). The
transverse stochastic motion is the implicit regulariser; smoothing the flow
removes exactly what makes the run generalise.

\subsection{The flip hazard field}
\label{sec:hazard}

Sign flips are governed by a single variable already present in the optimiser
state: $r_i = |m_i|/\sqrt{v_i}$, Adam's normalised update magnitude.

\begin{center}
\begin{tabular}{lrr}
\toprule
decile of $r$ & flip rate & vs.\ mean \\
\midrule
D1 (lowest)  & $32.0\%$ & $2.64\times$ \\
D5           & $\phantom{0}9.2\%$ & $0.76\times$ \\
D9           & $\phantom{0}0.8\%$ & $0.06\times$ \\
D10 (highest)& $\phantom{0}0.2\%$ & $0.01\times$ \\
\bottomrule
\end{tabular}
\end{center}

The hazard is monotone across every decile, spans $160\times$, and has AUC
$0.809$ as a single-variable predictor. The curve is not exponential but a
\emph{hard threshold}: above the $\sim\!70$th percentile of $r$, flips vanish
entirely (zero flips in the top three bins over $300$k sampled coordinates).
The flip mass concentrates where the normalised momentum is small, so the
dominant mechanism is low signal-to-noise near zero, not persistent opposing
gradients overwhelming a large momentum.

\begin{observation}[Aggregation destroys the signal]
The AUC $0.809$ is per-coordinate. Aggregated to tiles it collapses to
$0.60$--$0.64$, and lead-time prediction is flat (AUC $0.61$ at $k=1$ and at
$k=20$) with tile activity autocorrelation $\sim\!0.66$ out to $40$ steps and no
decay. Tiles are static ``hot regions,'' not dynamic boundary-approach regions.
Near-boundary is a per-coordinate property; the tile mean is dominated by the
frozen majority. The adaptive tile scheduler this was meant to enable does not
work.
\end{observation}

\subsection{Flips cluster, by neuron and by index}
\label{sec:cluster}

Per-tile flip fractions are over-dispersed relative to binomial independence.
A cover comparison at matched tile count and size separates the causes:

\begin{center}
\begin{tabular}{lr}
\toprule
cover & over-dispersion vs.\ binomial \\
\midrule
random             & $\phantom{1}1.00\times$ \\
input axis (column)& $1.36$--$1.57\times$ \\
output axis (row)  & $1.95$--$2.58\times$ \\
contiguous (large tiles) & $12.35\times$ \\
\bottomrule
\end{tabular}
\end{center}

Random calibrates the null at exactly $1.00\times$. Output-axis (neuron)
clustering exceeds input-axis by $\sim\!1.7\times$ in both a feed-forward and an
attention matrix, so flips cluster by the neuron a weight feeds, a functional
rather than storage grouping. The large $12.35\times$ for $4228$-parameter tiles
is mostly cross-tensor heterogeneity; within a single matrix the effect is
$2$--$3\times$. Strided and frieze covers give $\sim\!3.5\times$, so roughly half
the clustering is global-in-index and half is adjacency.

\subsection{A periodic cover gives the first non-forced obstruction}
\label{sec:frieze}

Every earlier cohomology vanished by construction: contractible bases, path
nerves, an AF Bratteli diagram. A frieze (periodic) cover makes the nerve a
\emph{cycle}, which need not be trivial. Fitting each of the $256$ edge maps
independently and transporting around the loop:

\begin{center}
\begin{tabular}{lrrr}
\toprule
holonomy $H$ & $\|H-I\|$ & eigenvalues $|\cdot|$ & $\det$ \\
\midrule
shared $M^{256}$        & $0.4498$ & $(1.000,\,0.636)$ & $0.636$ \\
per-edge $\prod_i M_i$  & $0.3424$ & $(1.000,\,0.729)$ & $0.729$ \\
\bottomrule
\end{tabular}
\end{center}

\begin{proposition}[Genuine holonomy with a preserved line]
Independently-fitted edge maps accumulate a non-identity holonomy around the
cycle, so it is not a shared-fit artefact. Each $M_i$ is near-unimodular
($\overline{|\det|}=0.999$, spread $1.7\%$) yet the product has $\det=0.73$: the
contraction accumulates from many tiny consistent deviations, the signature of a
connection with small constant curvature. The loop map fixes one direction
(eigenvalue $1.000$) and contracts the transverse one ($0.73$), so
$H^0\cong\ker(H-I)$ is one-dimensional --- global sections exist in one stalk
coordinate and fail in the other. This is the first cohomological obstruction in
the programme that reports the data rather than the graph.
\end{proposition}

The curvature is real but localised: per-plaquette, $88.2\%$ have
$|F_\square|<0.05$, and curvature magnitude correlates with the variation of $r$
around the plaquette (Spearman $0.590$; $8.4\times$ split between low- and
high-$\nabla r$ halves). The connection is flat almost everywhere with defects
at hazard-field gradients.

\subsection{The sign field is compressible; adjacency is nearly redundant}
\label{sec:signentropy}

The conditional entropy of the update sign, in bits per parameter per step:

\begin{center}
\begin{tabular}{lrr}
\toprule
conditioning & bits/param/step & vs.\ raw \\
\midrule
none (raw sign field) & $1.000$ & $1.000$ \\
previous step's sign & $0.557$ & $0.557$ \\
$+$ neuron (row) context & $0.541$ & $0.541$ \\
$+$ Adam $r=|m|/\sqrt{v}$ & $0.428$ & $0.428$ \\
$+$ $r$ and neuron context & $0.405$ & $0.405$ \\
\bottomrule
\end{tabular}
\end{center}

The sign field carries $0.41$ bits per parameter per step, not $1$: a
$79\times$ compression against fp32 when combined with the single magnitude
scalar. Temporal context does most of the work; $r$ adds $23\%$ and is free,
already in the optimiser state; neuron adjacency adds only $3$--$5\%$. Despite
the neuron clustering of \S\ref{sec:cluster}, the sign domains are essentially
$r$-level sets, so adjacency carries little information beyond the hazard
variable. This is a bandwidth result, not a compute result: the backward pass
is still needed to learn the signs.

\subsection{The residual: Fisher-large, trajectory-null}
\label{sec:residual}

Does $u\approx a\cdot\operatorname{sign}(u)$ discard anything that matters? Under
diagonal Fisher the reconstruction retains $R=0.64$, barely improving with
$16{,}384\times$ more magnitude parameters ($0.636\to0.697$), while training
retains $100.1\%$. The metric escape --- that diagonal Fisher underweights the
residual --- fails in the informative direction:

\begin{center}
\begin{tabular}{lr}
\toprule
metric & retention of $a\cdot\operatorname{sign}(u)$ \\
\midrule
diagonal Fisher & $0.630 \pm 0.020$ \\
full Fisher (finite-difference JVP) & $0.417 \pm 0.078$ \\
\bottomrule
\end{tabular}
\end{center}

The true Fisher retains \emph{less} ($0.42$), so the discarded residual
$\varepsilon=u-a\operatorname{sign}(u)$ carries $58\%$ of the true Fisher
energy, concentrated in the correlated directions the diagonal cannot see ---
and training is unaffected. The resolution is geometric.

\begin{proposition}[The residual is transverse to descent]
$\cos^2(\varepsilon,g)=0.031$ (Euclidean) and $0.034$ (Fisher): the residual is
$97\%$ perpendicular to the gradient. Fisher energy measures
$\varepsilon^\top G\varepsilon$ and is large; first-order loss change measures
$\langle g,\varepsilon\rangle$ and is near zero, because $\varepsilon$ points
across the level set rather than down it. The small parallel part
($\langle g,\varepsilon\rangle/\langle g,u\rangle=+0.38$) mildly \emph{opposes}
descent, so $a\operatorname{sign}(u)$ retains $62\%$ of the descent pairing while
shedding a slightly counterproductive residual. ``$58\%$ of Fisher energy
discarded'' and ``training unchanged'' are therefore not in tension: they are
the two faces of a transverse vector.
\end{proposition}

The residual is structured, not noise: temporal autocorrelation $0.62$ at lag
one, decaying to noise by lag ten; it does not predict future gradients
($\approx 0$ at all future lags); and its magnitude correlates with $r$ at
$0.88$. The energy is bimodal in $r$ --- $45\%$ in the top decile (frozen
coordinates, where one magnitude misfits a near-zero distribution) and $24\%$ in
the bottom two (active coordinates, where the flipping is itself the transverse
motion), with a near-empty middle. The strong claim ``$a\cdot\operatorname{sign}$
is the descent projection'' is \emph{not} supported: both $u$ and its
reconstruction are $\sim\!95\%$ transverse, and the sign step simply keeps enough
of the descent sliver while the residual keeps almost none. What is established
is weaker and more robust: Fisher-norm importance and gradient-alignment
importance are different notions of importance, and the residual is large in the
first and null in the second. Whether accumulated transverse motion has
second-order consequences over long horizons is not settled by these
instantaneous measurements.

\subsection{Progression, and the object that survived}
\label{sec:progression}

The programme began seeking a reduced observable from which weights could be
computed without gradient descent. It ends with that goal refuted on two
independent grounds and a narrower object in hand. The honest sequence:

\emph{Refuted, algebraic closures.} The symmetric flow product does not close
($0.7\%$, \S\ref{sec:phase3-algebra}); no $A_\infty$ correction exists (Stasheff
residual $1.0$); the Lie structure is local, decaying $30\times$ over $40$ steps
(\S\ref{sec:lie}); the tile data is not a vector bundle (cocycle fails
$5$--$80\%$, plaquette holonomy $0.04$--$1.29$); refinement and the Frobenius
relation do not commute. Each was tested and failed, which is what excludes
treating the surviving object as a static algebraic structure.

\emph{Refuted, the reduced-observable programme.} $\beta$ is cover-dependent
($\sim\!9\%$ of per-parameter variance, \S\ref{sec:cover-dependence}) and carries
no directional content (\S\ref{sec:magdir}); the strand ultrametric is
contiguity (permuted tree $\to 0.000$, \S\ref{sec:strands}); the topological
invariants vanish by construction ($K_1=0$ by Elliott, the index by dimension
counting, $H^\bullet$ by contractible bases). The one closed-loop holonomy test
in parameter space returned trivial ($w_1=0$, every block $+1$).

\emph{Established, and reproducible.} The driftless-walk law and $E(64)$
($\sim\!50\sigma$, three seeds); the rank-one curvature response, which derives
the scalar-per-block preconditioner; the Fredholm cokernel ($59$--$72\%$, flat,
largest at the minimum), which explains why block-local methods cannot close the
gap; the LayerNorm spectral gap, surviving occupancy and weight/bias controls;
and the frieze holonomy (\S\ref{sec:frieze}), the sole non-forced obstruction.

\emph{The surviving object.} It is not an algebra. It is a controlled
stochastic process on the pair $(r,s)$, where $r=|m|/\sqrt v$ is a scalar
stability field and $s\in\{\pm1\}^N$ the sign field, with the update
reconstructed as $u_t\approx a_t\, s_t$. The dynamics live on $(r,s)$, not on
$u$. Five independently measured phenomena organise under $r$: flip hazard
(AUC $0.809$, \S\ref{sec:hazard}), sign entropy ($23\%$ reduction,
\S\ref{sec:signentropy}), plaquette curvature (localised at $\nabla r$,
\S\ref{sec:frieze}), coordinate freezing (smooth saturation of the gap
$1-2h(r)$), and residual magnitude (correlation $0.88$, \S\ref{sec:residual}).
Every attempt to give the object algebraic closure was measured and failed
(Frobenius $0.7\%$, Stasheff residual $1.0$, Lie decaying $30\times$, cocycle
$5$--$80\%$, AU coherence); it is dynamical, not algebraic.

\subsection{Controlled Sign Transport: the theory, and its regime}
\label{sec:cst}

The measurements assemble into a single structure with a bounded small
parameter. The transition law of the sign process factorises to leading order:
\[
  P(s_{t+1}\mid s_t, r_t) = K_{r_t}, \qquad
  K_r = K^{(0)}_r + \varepsilon\, K^{(1)}_r,
\]
where $K^{(0)}_r = \bigotimes_i \left(\begin{smallmatrix} 1-h(r_i) & h(r_i) \\ h(r_i) & 1-h(r_i)\end{smallmatrix}\right)$
is a product of independent two-state chains with hazard $h(r_i)$, and
$K^{(1)}_r$ is a weak neuron-supported coupling. Each layer of this
decomposition was reduced by a control, not proposed by analogy.

\emph{The backbone $K^{(0)}$ is trivial and carries most of the behaviour.} Its
spectral gap is $1-2h(r)$, which is the hazard curve rewritten, so no index-like
transition beyond the freezing crossover exists there. The gap rises smoothly
from $0.175$ (lowest $r$-octile) to $0.974$ (highest); freezing is a saturation,
not a phase transition, and no critical $r_c$ appears.

\emph{The coupling $K^{(1)}$ is the only non-factorised content.} Conditioning on
both coordinates' $r$, adjacent (same-neuron) coordinates still flip together
$27\%$ more than independence predicts, five times the residual of random distant
pairs. Spectrally, $11$ of $128$ flip-covariance eigenvalues exceed the
shuffled-independent null by $3\sigma$ (top mode $10.9\sigma$), and their
magnitude ($23\%$ excess) matches the pair coupling: the collective modes
\emph{are} the neuron coupling. The interesting operator theory is
$\sigma(K^{(1)})$, not $\lambda_2(K)$.

\emph{The holonomy is abelian.} The frieze transport maps commute to four
decimals ($\|[M_i,M_j]\|/\|M_i\|\|M_j\| < 10^{-3}$), and the ordered loop product
has the same eigenvalues as the order-free product of eigenvalues
($(1.000,0.729)$ vs $(1.004,0.726)$). The holonomy is therefore a rank-one
dissipative phase --- one fixed direction, one contraction to $0.73$, one scalar
per loop --- and $F = dA$ with $A\wedge A \approx 0$: curvature comes from
spatial variation of the scalar control, not from non-commutativity. Geometric
language is available but optional; the object is one number attached to a loop.

\begin{proposition}[The perturbative regime is the computational parameters]
The coupling $\varepsilon$ (relative joint-flip excess after conditioning on $r$)
is small and stable across the transformer blocks and large in the embedding:
\end{proposition}

\begin{center}
\begin{tabular}{lrr}
\toprule
tensor & $\varepsilon$ early (steps 1--100) & $\varepsilon$ late (101--200) \\
\midrule
FF gate, $L_2$      & $24\%$ & $25\%$ \\
FF out, $L_2$       & $11\%$ & $\phantom{0}6\%$ \\
$W_Q$, $L_2$        & $24\%$ & $27\%$ \\
$W_V$, $L_2$        & $14\%$ & $12\%$ \\
$W_O$, $L_2$        & $14\%$ & $12\%$ \\
FF gate, $L_0$      & $24\%$ & $27\%$ \\
FF gate, $L_5$      & $27\%$ & $31\%$ \\
embedding           & $211\%$ & $267\%$ \\
\bottomrule
\end{tabular}
\end{center}

For the block parameters --- the parameters implementing the function ---
$\varepsilon \in [0.06, 0.31]$, flat between early and late training, so the
decomposition is genuinely perturbative and the sign process is dominated by
$r$-controlled independent flips. For the embedding --- the lookup table storing
the single repeated corpus --- $\varepsilon > 2$ and grows with training, so the
coupling dominates and the independent-flip backbone does not describe the
dynamics. The boundary was not imposed; it fell out of the sweep, and it
separates distributed computation (perturbative) from associative memory
(non-perturbative). Controlled Sign Transport holds in the former regime.

\emph{Generation versus transport.} The representation $(a,r,s)$ reconstructs the
update but cannot be produced without the backward pass: $r$ and $s$ are outputs
of backprop, not substitutes for it. This is the terminal statement of the
programme. Computation separates into generation --- which requires the
gradient and remains the irreducible cost --- and transport --- which requires
only $(a,r,s)$ and is therefore cheap to store, transmit, and partially reuse.
The two levers that survive a compute or bandwidth accounting are transport-side:
stale-sign at matched backward passes (\S\ref{sec:stalesign}) and the $79\times$
sign-field compression (\S\ref{sec:signentropy}). Neither reduces the cost of
generation. The programme began seeking to replace gradient descent and ends
having characterised, rather than removed, the one operation that cannot be
replaced.

\begin{figure*}[t]
\centering
\resizebox{\textwidth}{!}{%
\begin{tikzpicture}[
  font=\small,
  box/.style={draw, rounded corners=1pt, align=center, inner sep=4pt, minimum height=7mm},
  disc/.style={draw, double, rounded corners=1pt, align=center, inner sep=4pt, fill=black!4},
  strand/.style={draw, align=center, inner sep=3pt},
  small/.style={draw, rounded corners=1pt, align=center, inner sep=2pt, font=\scriptsize},
  flow/.style={-{Latex[length=2mm]}, thick},
  el/.style={font=\footnotesize\itshape, text=black!60}
]
% ===== LEFT COLUMN: pipeline =====
\node (corpus) [box] {CORPUS \\ \footnotesize batch / tokens};
\node (force)  [box, below=4mm of corpus] {FORCING FIELD \\ {\footnotesize\itshape what is presented}};
\node (discL)  [disc, below=4mm of force, text width=52mm]
  {DISC\; (slow variables) \\ {\footnotesize\itshape geometry of possible motion} \\[1pt] $r,\;m,\;v$ history coupling};
\node (bwd)    [below=4mm of discL, align=center] {\textbf{BACKWARD PASS} \\ \footnotesize measurement (irreducible)};
\node (strand) [strand, below=4mm of bwd, text width=56mm]
  {\textbf{STRAND}\;\; $S_t=\operatorname{sign}(\text{gradient})$ \\[3pt]
   \begin{tabular}{@{}p{22mm}p{26mm}@{}}
   \scriptsize persistent core & \scriptsize moving frontier\\
   \tiny$\blacksquare\blacksquare\blacksquare\blacksquare\blacksquare\blacksquare$ &
   \tiny$\square\,\cdot\,\square\;\cdot\,\square\,\cdot$\\
   \tiny$\blacksquare\blacksquare\blacksquare\blacksquare\blacksquare\blacksquare$ &
   \tiny$\cdot\,\square\,\cdot\;\square\,\cdot\,\square$\\
   \end{tabular}};
\node (transL) [box, below=5mm of strand, text width=50mm]
  {TRANSPORT \\ \footnotesize stale sign \;\textbullet\; entropy compression \;\textbullet\; communication};
\node (theta)  [below=3mm of transL] {$\theta(t{+}1)$ \;\footnotesize next state};
\draw[flow] (corpus)--(force);
\draw[flow] (force)--(discL);
\draw[flow] (discL)--(bwd);
\draw[flow] (bwd)--(strand);
\draw[flow] (strand.240) to[out=-90,in=90] node[left,el,pos=0.9]{stable} (transL.140);
\draw[flow] (strand.300) to[out=-90,in=90] node[right,el,pos=0.9]{churn} (transL.40);
\draw[flow] (transL)--(theta);

% ===== MIDDLE COLUMN: the dynamic loop =====
\begin{scope}[shift={(66mm,0)}]
  \node (l0) {\footnotesize\textsc{dynamic loop}};
  \node (l1) [below=3mm of l0] {next batch};
  \node (l2) [below=3mm of l1] {corpus};
  \node (l3) [below=3mm of l2] {perturbation};
  \node (l4) [below=3mm of l3] {backward};
  \node (l5) [below=3mm of l4] {strand};
  \node (l6) [below=3mm of l5] {parameter update};
  \node (l7) [below=3mm of l6] {new geometry};
  \foreach \i/\j in {l1/l2,l2/l3,l3/l4,l4/l5,l5/l6,l6/l7}
    \draw[flow] (\i)--(\j);
  \draw[flow] (l7.east) to[out=0,in=0,looseness=1.6] (l1.east);
\end{scope}

% ===== RIGHT COLUMN: disc manifold + strand fields =====
\begin{scope}[shift={(112mm,-4mm)}]
  \node (ef) [align=center] {\footnotesize external forcing};
  \draw[thick, fill=black!5]
    (-16mm,-6mm) .. controls (-6mm,4mm) and (6mm,-2mm) .. (16mm,4mm)
    -- (18mm,-11mm) .. controls (6mm,-9mm) and (-6mm,-15mm) .. (-14mm,-17mm) -- cycle;
  \node[align=center] at (1mm,-5mm) {DISC \\ \scriptsize high-dim.\ smooth};
  \node[font=\scriptsize, align=center] at (1mm,-14mm) {$\rho(\text{batch dist},\text{sim}){=}0.62$};
  \foreach \d in {2,1,0}
    \node[strand, fill=white, text width=40mm, minimum height=12mm] (sf) at ($(0,-32mm)+(\d mm,\d mm)$) {};
  \node[align=center] at (1mm,-29mm) {\footnotesize STRAND FIELD};
  \node[font=\tiny] at (1mm,-34mm) {$\blacksquare$ core \quad $\square$ frontier};
  \draw[flow] (ef) -- (1mm,-3mm);
  \draw[flow] (1mm,-16mm) -- (1mm,-25mm);
\end{scope}

% ===== BOTTOM ROW: the batch->strand manifold result =====
\node (bes) [box, below=13mm of transL, text width=46mm, xshift=6mm]
  {BATCH EMBEDDING SPACE \\ {\footnotesize\itshape smooth, high-rank map}};
\node (lin) [small, left=10mm of bes]  {linear \\ $0.699$};
\node (knn) [small, above left=4mm and 10mm of bes] {kNN $k{=}15$ \\ $0.704$ (plateaus)};
\node (nl)  [small, right=10mm of bes] {nonlinear (rand.\ feat.) \\ $0.535$ \;\textbf{overfits}};
\node (verdict) [below=4mm of bes, text width=72mm, align=center, font=\footnotesize]
  {low-dimensional curved manifold \textbf{disproven} (trivial corpus): nonlinear
   does \emph{not} beat linear.\; Residual $\approx\!0.30$ is high-dimensional
   noise in $\delta$.\; \textsc{proven:} high-dimensional continuous field,
   compressible only in transport (linear baseline $0.699$).};
\draw[flow] (lin)--(bes); \draw[flow] (knn)--(bes); \draw[flow] (nl)--(bes);
\draw[flow] (bes)--(verdict);
\end{tikzpicture}%
}
\caption{The disc--strand picture as measured. \textbf{Left:} the pipeline ---
a smooth, high-dimensional corpus \emph{forcing field} selects the strand's
\emph{identity} (not its magnitude); the \textbf{disc} $(r,m,v)$ is the slow
manifold (geometry of possible motion, not a controller: it sets persistence and
the fixed $\sim\!55\%$/step churn velocity but cannot generate or gate the
strand); the irreducible \textbf{backward pass} produces
$S_t=\operatorname{sign}(g_t)$, split into a persistent core ($\sim\!115$k
coordinates) and a frontier sweeping $99.6\%$ of the network; only
\emph{transport} downstream compresses (stale-sign, $79\times$ sign entropy).
\textbf{Middle:} the training loop --- corpus perturbation, measurement, strand,
update, new geometry. \textbf{Right:} the disc as a high-dimensional smooth
manifold; nearby batches give nearby strands (batch-distance/similarity
$\rho{=}0.62$), strand fields orbiting it with the core/frontier split.
\textbf{Bottom:} the batch$\rightarrow$strand map is smooth but high-rank ---
linear predicts strand sign at $0.699$, kNN plateaus at $0.704$, and a nonlinear
model \emph{fails} to beat linear ($0.535$, overfits), disproving a
low-dimensional curved manifold for this corpus. The corpus signal is real and
continuous but high-dimensional, with a $\approx\!0.30$ residual living in the
backward error $\delta$; whether a structured corpus yields a low-dimensional map
is the one open question (\S\ref{sec:backward}).}
\label{fig:disc-strand}
\end{figure*}

\subsection{The backward pass: what is irreducible, and why}
\label{sec:backward}

The generation--transport split of \S\ref{sec:cst} leaves one question: can the
generation step --- the backward pass --- itself be made cheaper? Eight
experiments attack it from independent directions and return one wall.

\emph{Prediction and pruning fail.} The next active frontier (coordinates whose
sign is about to change) is predictable from pre-backward optimiser state only
to AUC-equivalent recall $0.62$ at $20\%$ selection, and a linear model on
$(r,\Delta r,|m|,v)$ adds nothing over $r$ alone. Adding the current batch's
forward activations --- the most direct encoding of the incoming sample ---
leaves recall at chance for the batch feature, because a weight's gradient is
$\delta_{\text{out}}\!\cdot\!\text{act}_{\text{in}}$ and the forward pass
supplies only the second factor; the first is the backward pass itself.

\emph{Sketching fails on structured error.} The update sign is invariant to
i.i.d.\ gradient noise up to $25\%$ ($100\%$ agreement) but breaks under
rank-$25\%$ SVD truncation ($82\%$ agreement, $1.6\times$ val cost) and under any
low-rank surrogate. The enemy is correlated error, which is exactly what every
compute-saving spatial approximation introduces. Neither a per-neuron
mean-removal nor the graph-Laplacian basis of the flip-coupling graph recovers
the signs, because that graph is dense and non-modular (degree $140$ of $512$,
no spectral gap): there is no privileged basis to find.

\emph{Gating fails worse than random.} Training under a rule that computes the
exact gradient on a selected $20\%$ and persists the strand elsewhere:

\begin{center}
\begin{tabular}{lrr}
\toprule
gate & granularity & val vs.\ full \\
\midrule
random    & coordinate & $3.4\times$ \\
$r$-gated & coordinate & $6.9\times$ \\
random    & neuron-tile & $3.4\times$ \\
$r$-gated & neuron-tile & $23\times$ \\
momentum  & neuron-tile & $24\times$ \\
\bottomrule
\end{tabular}
\end{center}

\begin{proposition}[Coupling-aligned gating is anti-optimal]
Random selection is the floor; every principled gate does strictly worse, and
the damage scales with how tightly the gate concentrates the computed/persisted
boundary onto the coupling structure. Neuron-tile gating, which aligns the cut
with the between-neuron couplings, is worse than coordinate gating; both are
worse than random, which has zero coupling-alignment. A coordinate's importance
is non-local --- set by its coupling to the coordinates one would skip --- so
any selection rule correlated with that coupling cuts precisely what must be
preserved.
\end{proposition}

The mechanism is the absence of scale separation. Multigrid and adaptive-mesh
methods reduce cost when error is smooth with local corrections; the sign field
is a densely coupled field with no hierarchy (no modularity, no spectral gap, no
low-rank structure), which is exactly the regime where hierarchical and sparse
methods provably do not apply.

\begin{observation}[Terminal statement]
The strand --- a temporally coherent, $r$-organised sign field --- is the
irreducible object of the optimisation. Backpropagation is the only known
mechanism that generates it. Its cost cannot be reduced by sparsity in any form
(coordinate, tile, predicted, or sketched), because importance is non-local:
dense, non-hierarchical coupling without scale separation. What remains
compressible is transport, not generation: stale-sign across steps
(\S\ref{sec:stalesign}) and sign entropy across coordinates
(\S\ref{sec:signentropy}), both operating on the already-measured field. Whether
a different \emph{dense} mechanism --- feedback alignment through the recursion,
forward-gradient, or a synthetic backward field --- can generate the strand more
cheaply while preserving weighted sign agreement above $\sim\!0.95$ through
network depth under correlated error is the single question this programme
leaves open. Sparsity-based economy is refuted; substitution is untested.
\end{observation}

The programme set out to ask whether the disc--strand decomposition could
replace gradient descent. The answer is more precise than a refusal: disc and
strand are a complete, compressible, transportable description of what gradient
descent \emph{produces}, and the reason they cannot replace what produces it is
a measurable structural fact about the loss geometry --- non-locality without
scale separation --- rather than a limitation of the representation.

\section{Phase 4: Corpus Generality of the High-Dimensionality Result}
\label{sec:corpus-generality}

Every measurement to \S\ref{sec:backward} was taken on one corpus: a single
sentence looped roughly $300$ times. The central negative finding --- that the
strand field is high-dimensional and therefore does not compress across corpus
states --- was therefore open to the objection that it is an artefact of that
triviality. A structured corpus might make the batch~$\to$~strand map
low-dimensional, which would reopen dictionary compression and make the
sparsity refutation corpus-specific rather than general. This section settles
that objection.

\subsection{The structure axis and its confound}
\label{sec:axis}

Six corpora at $V=256$, $60{,}000$ tokens each, spanning a structure axis
validated two ways: gzip compression ratio, and the block-entropy gap at
$L=6$ against each corpus's own token shuffle. The shuffle controls are the
only confound-free comparison available, since they hold vocabulary and
unigram statistics fixed and vary order alone.

\begin{center}
\begin{tabular}{lrrr}
\hline
corpus & live vocab & gzip ratio & blockH gap \\
\hline
\texttt{iid}          & $256$ & $1.001$ & $0.000$ \\
\texttt{english\_shuf}& $27$  & $0.631$ & $-0.000$ \\
\texttt{english}      & $27$  & $0.282$ & $0.362$ \\
\texttt{code\_shuf}   & $95$  & $0.663$ & $-0.000$ \\
\texttt{code}         & $95$  & $0.259$ & $0.436$ \\
\texttt{repeat}       & $97$  & $0.007$ & $1.494$ \\
\hline
\end{tabular}
\end{center}

The live-vocabulary column carries a confound that governs how the results may
be read. Vocabulary size is not constant along the axis --- it runs
$256, 27, 95, 97$ --- so the cross-corpus ordering conflates structure with
alphabet size and unigram entropy. Only the paired contrasts
(\texttt{english} versus \texttt{english\_shuf}, \texttt{code} versus
\texttt{code\_shuf}) isolate structure. The cross-corpus trend is reported
below but is not load-bearing.

\subsection{Protocol, and two corrections to it}
\label{sec:idprobe-protocol}

For each corpus: train to a checkpoint at $80$ steps, freeze $\theta$, sweep
$N_B$ batches, and record $s_B=\operatorname{sign}(g_B)$ on a fixed random
subset of $20{,}000$ block-parameter coordinates together with a batch
embedding $z_B$ (mean token embedding concatenated with the normalised token
histogram). Predict $s_B$ from $z_B$ linearly, nonlinearly by random cosine
features, and by $k$-nearest neighbours; report the participation ratio of the
strand spectrum and the correlation between batch distance and strand
dissimilarity.

Two corrections were necessary, and both change what the numbers mean.

\emph{First, the original comparison was between two interpolators.} With
$z_B$ of dimension $D+V=512$ and $N_B=200$ giving $150$ training rows, both the
linear least-squares fit and the $800$-dimensional random-feature fit are
underdetermined minimum-norm solutions. The earlier reading --- nonlinear
$0.535$, ``overfits,'' against linear $0.699$ --- is therefore substantially a
statement about minimum-norm behaviour at $p \gg n$, not about the intrinsic
dimension of the map. Here $N_B=600$, giving $360$ training rows, and both
model classes are placed on a ridge path with the penalty selected on a
held-out split, so each is given its best available fit. Untuned numbers are
reported alongside for continuity.

\emph{Second, the initialisation must not depend on the corpus.} The
compiler prefix constructs the embedding from the corpus bigram Laplacian via
$\operatorname{eigsh}(L_{\mathrm{sym}},k=D{+}1)$. At $V\le 256$ this is not
merely inadvisable but impossible, since $k=257$ exceeds the matrix dimension;
more importantly it would make $\theta_0$ a function of the corpus, and
differentially so --- \texttt{english} has $27$ live tokens against
\texttt{iid}'s $256$, so the graph is degenerate to different degrees along the
axis. A probe asking whether corpus structure lowers the dimension of the
batch~$\to$~strand map cannot have corpus structure pre-loaded into the
initialisation at corpus-dependent strength. All runs therefore use $V=256$
fixed and a plain seeded initialisation, identical across corpora. This
forfeits comparability with the historical $V=1017$ numbers, which is the
correct trade for a cross-corpus comparison and the wrong one for continuity
with earlier phases.

\subsection{Result: the wall holds}
\label{sec:idprobe-result}

Block parameters, $N_B=600$, split $360/90/150$, single seed.

\begin{center}
\begin{tabular}{lrrrrrrr}
\hline
corpus & lin (raw) & nl (raw) & lin (ridge) & nl (ridge) & gain & kNN$_{15}$ & PR \\
\hline
\texttt{iid}          & $0.587$ & $0.527$ & $0.729$ & $0.597$ & $-0.133$ & $0.646$ & $208.0$ \\
\texttt{english\_shuf}& $0.719$ & $0.579$ & $0.842$ & $0.663$ & $-0.180$ & $0.751$ & $69.9$ \\
\texttt{english}      & $0.748$ & $0.639$ & $0.849$ & $0.716$ & $-0.134$ & $0.787$ & $65.1$ \\
\texttt{code\_shuf}   & $0.670$ & $0.590$ & $0.817$ & $0.672$ & $-0.145$ & $0.736$ & $80.4$ \\
\texttt{code}         & $0.572$ & $0.557$ & $0.691$ & $0.652$ & $-0.039$ & $0.663$ & $243.6$ \\
\texttt{repeat}       & $0.745$ & $0.661$ & $0.865$ & $0.756$ & $-0.109$ & $0.769$ & $57.3$ \\
\hline
\end{tabular}
\end{center}

\begin{observation}[The discriminator returns \emph{no}]
\label{obs:no-manifold}
The nonlinear model does not beat the linear model on any corpus on the axis.
The gain is negative throughout, from $-0.039$ to $-0.180$, and $k$-nearest
neighbours --- which requires no regularisation choice and is therefore the
independent check --- tracks the linear model from below everywhere. Two
structurally different non-parametric estimators both fail to beat a linear
map. The high-dimensionality of the batch~$\to$~strand map is not an artefact
of the degenerate corpus, and dictionary compression does not reopen.
\end{observation}

Ridge tuning raised the linear model substantially ($0.748\to0.849$ on
\texttt{english}) and raised the nonlinear model too, without closing the gap:
the earlier ``nonlinear overfits'' reading was partly an artefact of the
minimum-norm fit, and the conclusion survives correcting it.

The paired contrasts are where caution is required.

\begin{center}
\begin{tabular}{lrrr}
\hline
pair & lin (ridge) & nl (ridge) $\to$ gain & PR \\
\hline
\texttt{english\_shuf}$\to$\texttt{english} & $0.842\to0.849$ & $-0.180\to-0.134$ & $69.9\to65.1$ \\
\texttt{code\_shuf}$\to$\texttt{code}       & $0.817\to0.691$ & $-0.145\to-0.039$ & $80.4\to243.6$ \\
\hline
\end{tabular}
\end{center}

Structure moves the gain toward compressibility in both pairs ($+0.046$,
$+0.106$), so this is not a flat null; but the sign never reverses, and the two
pairs disagree on the dimension measure. On \texttt{english} the participation
ratio falls slightly with structure, as the dictionary hypothesis predicts. On
\texttt{code} it triples, while linear predictability falls from $0.817$ to
$0.691$ and the distance--dissimilarity correlation falls from $0.561$ to
$0.318$. Real code structure makes the strand field \emph{less} predictable and
\emph{higher}-dimensional --- the opposite of a low-dimensional curved
manifold, and the single most informative number in the table.

\begin{observation}[Structured but not manifold-structured]
\label{obs:structured-not-manifold}
Participation ratios of $57$--$244$ against a ceiling of $600$ are well below
what independent sign vectors would give. The strand field carries real
redundancy, consistent with the $0.405$ bits/parameter/step of
\S\ref{sec:signentropy}; it is simply not organised as a low-dimensional curved
base. High-dimensional and compressible are not in tension --- this is the
transport/generation split of \S\ref{sec:backward} seen from the corpus side.
\end{observation}

\subsection{Replication}
\label{sec:idprobe-repl}

The four corpora entering the paired contrasts were rerun at a second seed.

\begin{center}
\begin{tabular}{lrrrr}
\hline
pair & $\Delta$gain (s17) & $\Delta$gain (s23) & $\Delta$PR (s17) & $\Delta$PR (s23) \\
\hline
\texttt{english\_shuf}$\to$\texttt{english} & $+0.046$ & $+0.016$ & $-4.8$ & $-6.3$ \\
\texttt{code\_shuf}$\to$\texttt{code}       & $+0.106$ & $+0.097$ & $+163.3$ & $+179.7$ \\
\hline
\end{tabular}
\end{center}

The sign of every gain replicates and no gain approaches zero, so
Observation~\ref{obs:no-manifold} is not seed-dependent. The two paired deltas
behave differently under replication: the \texttt{code} delta is stable
($+0.106$, $+0.097$) while the \texttt{english} delta is not ($+0.046$,
$+0.016$), so only the former should be treated as a measured effect. The
divergence of the two pairs on the dimension measure replicates cleanly and is
the more robust finding: \texttt{english} loses a few units of participation
ratio with structure, \texttt{code} gains over $160$.

\emph{Limitations.} Two seeds, one checkpoint at $80$ steps, one architecture. The negative result is a statement
about one nonlinear family --- $800$ random cosine features --- and the
agreement of $k$NN is what licenses the broader reading. The embedding arm was
measured separately and is reported nowhere in this section, because
$\texttt{head.weight}$ is tied to $\texttt{te.weight}$ and its gradient is
dominated by a nearly batch-independent softmax term: that arm measures weight
tying, not associative memory, and testing the $\varepsilon$-table confound
would require an untied model.

\section{Phase 5: Symmetry, and the Stabilizer of the Strand}
\label{sec:symmetry}

If no low-dimensional manifold organises the strand, a weaker and more robust
object may still exist: a group of transformations under which the strand is
approximately invariant. Invariance does not require low-dimensional
coordinates, and a stabilizer subgroup would induce equivalence classes on
batch space that are cheaper to specify than the strands themselves. This
section constructs the apparatus to test that and reports what it returns.

\subsection{Synthetic corpora with exact symmetries}
\label{sec:synth}

Symmetry can only be tested where it is exact by construction, so the corpora
are generated from an explicit parse tree over a byte-level expression
grammar,
\[
S\to E\,;\qquad E\to T \mid E{+}T \mid E{-}T \qquad T\to F\mid T{*}F
\qquad F\to V\mid D\mid (E)
\]
with $V\in\{a..h\}$ and $D\in\{0..9\}$. Three exact symmetries of the
generator are available, with deliberately different statistical profiles:

\begin{center}
\begin{tabular}{llll}
\hline
transform & parse tree & token statistics & surface altered \\
\hline
\texttt{rename} & identical & identical      & $85\%$ \\
\texttt{comm}   & mirrored  & $\approx$ identical & $78\%$ \\
\texttt{ws}     & identical & \emph{changed} (gzip $0.503\to0.336$) & $84\%$ \\
\hline
\end{tabular}
\end{center}

\texttt{rename} applies a bijection to the variable alphabet; \texttt{comm}
swaps the operands of the commutative operators; \texttt{ws} inserts
whitespace around binary operators, leaving the parse tree untouched while
altering the byte statistics substantially. Both \texttt{rename} and
\texttt{comm} preserve length exactly, so paired batches align token for
token. A structural ladder of known distance (one terminal, one operator, one
subtree, full resample) and two controls complete the suite: \texttt{shuf}
destroys structure while preserving the histogram, and \texttt{lexctl} carries
\texttt{rename}'s histogram with its structure destroyed.

\texttt{lexctl} is the control that makes the experiment mean anything. Without
it, statistical similarity is indistinguishable from symmetry.

\subsection{The stratification, and why the prediction failed}
\label{sec:rstrat}

The intended discriminator was stratification by the disc. Under the reading of
\S\ref{sec:hazard}, coordinates above the hazard threshold should have
disc-determined sign and therefore agree between \emph{any} two batches, while
coordinates below it should agree at chance; only the middle band is
batch-conditioned. A genuine symmetry invariance would then have to show its
excess agreement in the middle band, and excess in the top band would indicate
disc-determined agreement being misread as invariance.

Measured, the independent-batch floor is flat across $r$ deciles:
$0.587,\,0.591,\,0.588,\,0.585,\,0.582,\,0.591,\,0.595,\,0.606,\,0.606,\,0.610$.
The excess for every condition is largest in the \emph{top} band, which is the
signature the prediction identified as spurious.

The diagnostic explains this, and the explanation is a statement about regime
rather than a rescue. On block parameters at this checkpoint the disc field has
median $r=0.001$, $r_{90}=0.061$, $r_{99}=0.247$, maximum $0.879$, with
$r>0.5$ holding for $0.09\%$ of coordinates. Since $r\to1$ is the
consistent-gradient limit, essentially every coordinate is noise-dominated:
there is no deterministic regime for the top decile to occupy. The $p_{90}/p_{10}$
spread is $159\times$, consistent with the $160\times$ monotone range of
\S\ref{sec:hazard}, and agreement does rise monotonically with decile
($0.587\to0.610$).

\begin{observation}[Assumptions unsatisfied, not prediction falsified]
\label{obs:regime}
The experiment sampled one asymptotic regime of the theory. The two-phase
prediction presupposes a deterministic population that this operating point
does not contain, so it is untested rather than confirmed or refuted. What the
data do show is that
$\partial_r\,P\!\left(s_i^{(A)}=s_i^{(B)}\right)>0$ already holds \emph{inside}
the purely stochastic regime --- a weak monotone law, and the kind of
quantitative statement a theory of the disc ought eventually to derive.
\end{observation}

\subsection{Result: the stabilizer is trivial at this operating point}
\label{sec:stabilizer-result}

Twenty-four trials, each comparing a base batch against matched variants.

\begin{center}
\begin{tabular}{lrrrr}
\hline
condition & agreement & excess & positional identity & histogram cosine \\
\hline
\texttt{comm}   & $0.887$ & $+0.293$ & $0.444$ & $1.000$ \\
\texttt{pert1}  & $0.814$ & $+0.220$ & $0.901$ & $0.997$ \\
\texttt{rename} & $0.770$ & $+0.176$ & $0.778$ & $0.995$ \\
\texttt{pert2}  & $0.767$ & $+0.173$ & $0.907$ & $0.990$ \\
\texttt{lexctl} & $0.747$ & $+0.153$ & $0.074$ & $0.995$ \\
\texttt{pert3}  & $0.692$ & $+0.098$ & $0.150$ & $0.990$ \\
\texttt{pert4}  & $0.607$ & $+0.013$ & $0.140$ & $0.982$ \\
\texttt{indep}  & $0.594$ & $0$      & $0.133$ & $0.981$ \\
\texttt{ws}     & $0.506$ & $-0.088$ & $0.037$ & $0.361$ \\
\hline
\end{tabular}
\end{center}

Three readings, which agree with each other.

\begin{observation}[Alpha-renaming is not a symmetry of the strand]
\label{obs:lexctl}
\texttt{rename} preserves the parse tree exactly and changes $85\%$ of the
surface, giving agreement $0.770$. Its structure-destroyed control at the same
histogram gives $0.747$. The apparent invariance decomposes into $0.153$ of
histogram effect and $0.023$ of structure effect. What looked like structural
invariance is low-order statistical similarity, which shuffling preserves.
\end{observation}

\begin{observation}[The strand follows surface statistics, not syntax]
\label{obs:ws}
\texttt{ws} leaves the parse tree identical and alters only surface statistics.
Its agreement is $0.506$, which is $0.088$ \emph{below} the unrelated-batch
floor: a transformation that does not touch the grammar destroys strand
agreement more thoroughly than an unrelated batch does.
\end{observation}

\begin{observation}[The structural ladder is a surface-distance ladder]
\label{obs:ladder}
Across the nine conditions, agreement correlates $+0.66$ with positional
identity and $+0.66$ with histogram cosine. The monotone decay along the
perturbation ladder ($0.814,0.767,0.692,0.607$ against a floor of $0.594$)
coincides with positional identity falling from $0.90$ to $0.15$. Once
\texttt{comm} ($0.887$ at positional identity $0.444$, histogram cosine
$1.000$) and \texttt{ws} ($0.506$ at $0.037$, $0.361$) are included, the
pattern resolves without a residual structural term: histogram similarity
dominates and positional identity modulates.
\end{observation}

\subsection{Does the stabilizer grow with training? Inconclusive, and why}
\label{sec:stabilizer-evolution}

A stabilizer measured at one checkpoint is a snapshot of a possibly moving
object. The natural follow-up is to hold the symmetry suite fixed and sweep it
along a training trajectory, watching the gap
$\texttt{rename}-\texttt{lexctl}$ --- structural invariance net of histogram
--- as a function of step. One run, checkpoints at
$20,40,80,160,320,640,1280$, twelve trials each, measurement never calling an
optimiser step.

\begin{center}
\begin{tabular}{rrrrrrr}
\hline
step & val loss & $r_{90}$ & $\Pr[r>0.5]$ & \texttt{rename} & \texttt{lexctl} & gap \\
\hline
$20$   & $2.934$ & $0.187$ & $0.0053$ & $0.817$ & $0.775$ & $+0.042$ \\
$40$   & $2.926$ & $0.057$ & $0.0013$ & $0.790$ & $0.772$ & $+0.018$ \\
$80$   & $2.928$ & $0.040$ & $0.0004$ & $0.797$ & $0.775$ & $+0.022$ \\
$160$  & $2.917$ & $0.056$ & $0.0006$ & $0.743$ & $0.736$ & $+0.007$ \\
$320$  & $2.922$ & $0.148$ & $0.0072$ & $0.739$ & $0.725$ & $+0.014$ \\
$640$  & $2.913$ & $0.085$ & $0.0006$ & $0.749$ & $0.711$ & $+0.039$ \\
$1280$ & $2.917$ & $0.270$ & $0.0240$ & $0.762$ & $0.695$ & $+0.066$ \\
\hline
\end{tabular}
\end{center}

\begin{observation}[The sweep is inconclusive by construction]
\label{obs:sweep-null}
The unigram entropy of the corpus is $2.914$ nats and the validation loss sits
at $2.913$--$2.934$ for the entire run. The model learned the token frequencies
and nothing further; it never acquired the grammar. A measurement of whether
the stabilizer acquires grammatical invariance therefore had no opportunity to
register one, and the apparent upward drift of the gap
($+0.042\to+0.066$) sits within noise at twelve trials, where the per-condition
standard deviation is approximately $0.05$.
\end{observation}

The cause is a design fault in the grammar rather than a property of strand
generation: terminals are drawn uniformly, so the syntax carries almost no
predictive mass relative to the random leaves, and a model at the unigram
entropy is already optimal against everything except the operator skeleton. The
fix is a generator whose structure is load-bearing for the loss --- skewed
terminal distributions, or dependencies between leaves --- so that val loss
falls materially below unigram entropy and there exists a learned
representation for a stabilizer to be a stabilizer of. Until that is run, the
question of \S\ref{sec:stabilizer-programme} is open.

\emph{Scope.} All Phase~5 statements are conditioned on an operating point at
which the model has demonstrably not learned the grammar. The correct reading
of \S\ref{sec:stabilizer-result} is therefore not that the symmetry group of
the strand is trivial, but that \emph{at this operating point strand similarity
is governed by low-order statistical similarity rather than grammatical
equivalence}. Whether a model that had learned the grammar would exhibit
grammatical invariance is untested. The methodological result ---
that without a histogram control one will mistake statistical similarity for
symmetry --- is not conditioned on the operating point and is the more portable
finding.

\section{Toward a Theory: What the Measurements Constrain}
\label{sec:theory}

What follows is reformulation, not evidence. It reorganises measurements
already reported and adds no independent confirmation of any of them. Its
purpose is to record which mathematical descriptions are consistent with the
measurements, which are excluded by them, and --- for the parts that remain
conjectural --- what would decide them. Status is marked throughout.

\subsection{Two retractions}
\label{sec:retractions}

Two claims made in the course of this programme's discussion were stronger than
the evidence and are withdrawn.

\emph{``There is no geometry.''} Dense coupling without a spectral gap
(\S\ref{sec:backward}) excludes \emph{spatial locality in the parameter index}.
It does not exclude geometry. Spin glasses, kernel methods, Gaussian processes
and random matrix ensembles all combine dense interaction with rich spectral,
statistical or information-geometric structure. The defensible statement is the
narrow one: there is no useful metric on the parameter index.

\emph{``The correspondence with kinetic Ising is exact.''} Glauber dynamics is
defined by an energy function, and no Gibbs form for $K_r$ has been exhibited.
The dictionary --- strand to spin configuration, $r$ to quenched site disorder,
$\varepsilon$ to weak coupling, the trivial product backbone to the
non-interacting limit --- is an \emph{effective model}, useful for the
predictions it generates and not an identification. Demonstrating that the
measured transition probabilities admit an effective Hamiltonian is a specific
open task, not a formality.

\subsection{The state space}
\label{sec:statespace}

The measurements do not support a geometry on parameters alone. The strand
depends jointly on model state, optimiser state and batch, so the object to be
given structure is
\[
\mathcal{X}\times\Theta\times\mathcal{O}
\;\longrightarrow\;\{\pm1\}^P ,
\]
with $\mathcal{X}$ the batch space, $\Theta$ the parameters and $\mathcal{O}$
the optimiser moments. In this factorisation the disc is a slowly varying field
over $\Theta\times\mathcal{O}$, the corpus selects a trajectory through
$\mathcal{X}$, and the strand is the response field. This accommodates directly
what \S\ref{sec:corpus-generality} and \S\ref{sec:symmetry} measured: corpus
identity strongly conditions strand identity, while the optimiser state governs
churn.

\subsection{The parameter index is a partition with a gauge group}
\label{sec:gauge}

\emph{Status: argued from architecture, with one testable consequence
outstanding.} The architecture carries exact permutation symmetries --- the
feed-forward hidden units, the attention heads, and the residual-stream
coordinates may be relabelled with matched relabelling of the matrices that
read and write them, leaving the represented function unchanged, and LayerNorm
does not break this. Any statement of the form ``index $i$ neighbours index
$j$'' is therefore a property of the labelling and not of the network, and no
observable may depend on it.

The sign map is equivariant under the group $\mathcal{D}$ generated by positive
diagonal rescalings and signed permutations, and under no larger subgroup of
$GL(P)$. This gives a compact reading of three results already recorded: sign
survives $25\%$ i.i.d.\ noise (a within-chart perturbation) while breaking under
rank truncation (which leaves $\mathcal{D}$), and the permuted refinement
ladder collapses to $0.000$ (\S\ref{sec:strands}) because the permutation was
not a signed permutation of the disc chart.

What replaces the lattice is a set with a nested architectural partition ---
network, layer, matrix, row, parameter --- carrying a permutation gauge group
acting within each cell. Nesting is gauge-invariant; adjacency within a level is
not. The consequent reinterpretation of the contiguity result is that
contiguity in the flat parameter vector is largely co-membership in a matrix
and row, so that the $0.000$ of the permuted ladder records the gauge symmetry
appearing as a null rather than the destruction of a metric. \emph{This is a
prediction and it is untested}: conditioning on same-matrix and same-row
membership should remove the residual contiguity effect. If it does not, the
reinterpretation is wrong and there is genuine within-row geometry.

\subsection{Disc Riemannian, strand Finsler}
\label{sec:finsler}

\emph{Status: derivation from the measured update law; consistent with all
reported measurements, independently confirmed by none.} Steepest descent under
a norm constraint,
$u=\arg\min_{\|u\|\le a}\langle g,u\rangle$, gives $u\propto -g$ under
$\ell^2$, $u\propto -G^{-1}g$ under a Riemannian metric $G$, and
$u_i=-a\operatorname{sign}(g_i)$ under $\ell^\infty$, whose unit ball is the
hypercube and whose extreme points are exactly the sign vectors. The measured
law $u_t\approx a_t s_t$ of \S\ref{sec:sign-magnitude} is the third of these.

This is the reason a Riemannian formulation of the \emph{strand} cannot be
correct: a metric predicts a dense real-valued natural-gradient update, which is
measurably not what the optimiser produces. It also explains why the strand is
not a tensor --- gradients are covectors, and
$\operatorname{sign}((D\varphi)^{-T}g)\ne\operatorname{sign}(g)$ in general, so
the object that survived the programme is precisely the object that does not
transform covariantly.

The Fisher metric belongs instead to the disc, where it is already present
implicitly: $v_i$ estimates the diagonal empirical Fisher $\mathbb{E}_B[g_i^2]$,
so $G_\theta=\operatorname{diag}(\sqrt{v}+\epsilon)$ is a metric on $\Theta$
and $r_i=|m_i|/\sqrt{v_i}$ is the drift-to-noise ratio of coordinate $i$
measured in it. The hazard law of \S\ref{sec:hazard} then reads as a statement
about signal-to-noise: flips occur where Fisher-normalised drift is dominated
by Fisher-normalised noise, which predicts both the monotone flip rate and a
threshold. The disc is not a companion to the strand but the structure that
makes the strand well defined: $\mathcal{D}$ is the group preserving diagonal
Fisher structure, and without it there is no sign field to speak of.

\subsection{Frameworks excluded, and by which measurement}
\label{sec:excluded}

\emph{Status: exclusions, each traced to a specific measurement.} Recording
these prevents their recurrence.

\begin{center}
\begin{tabular}{lll}
\hline
framework & requires & excluded by \\
\hline
multigrid / refinement transport & scale separation & no spectral gap (\S\ref{sec:backward}) \\
persistent homology & filtration with separated scales & no spectral gap \\
lattice field theory & metric on the index & permutation gauge (\S\ref{sec:gauge}) \\
trajectory homology & non-contractible complex & trajectory is a path \\
index theory & stability under compact perturbation & rank truncation breaks sign \\
$K$-theory & additive or exact category & free category on a path \\
\hline
\end{tabular}
\end{center}

Three of these fail for one reason. The index-theoretic case is worth stating
separately: an index is invariant under compact perturbation, rank truncation is
finite-rank hence compact, and \S\ref{sec:backward} measured that the sign field
breaks under rank truncation while surviving i.i.d.\ noise --- so the
persistent/frontier decomposition has the opposite stability property to an
index. On $K$-theory, strands do not compose within strands, so the trajectory
generates a free category with no additive structure and $K_0$ is undefined;
passing to representations of the trajectory quiver gives
$K_0(\operatorname{rep}A_n)\cong\mathbb{Z}^n$, which counts steps. The honest
position is not that $K$-theory is dead but that no canonical object has
appeared to which it applies.

The cohomological content that does exist is already computed. The sign process
lives on $\{\pm1\}^P=(\mathbb{Z}/2)^P$, which is abelian, so abelian holonomy
is forced; the natural home for the non-forced part is
$H^1(\cdot\,;\mathbb{Z}/2)$, and the frieze cover of \S\ref{sec:frieze} is the
one non-forced class in it.

\subsection{The stabilizer programme}
\label{sec:stabilizer-programme}

\emph{Status: conjectural, and the one experiment attempted was inconclusive
(\S\ref{sec:stabilizer-evolution}).}

Phase~5 sought a fixed symmetry group of the data generator and found none.
The more defensible object is not fixed. The strand is a function of gradients,
gradients are functions of activations, and activations depend on whatever
statistics the representation has actually learned; so any invariance of the
strand is an invariance of the \emph{learned representation}, not of the
generator. Early in training that is token statistics, which is exactly what
\S\ref{sec:stabilizer-result} measured. Whether it later becomes syntax is
undetermined.

This reframes the question. Rather than asking which transformations preserve
the strand, ask how the stabilizer subgroup
\[
\operatorname{Stab}_t=\{\sigma:\ s_{\sigma(B)}\approx s_B \text{ at }
(\theta_t,\mathcal{O}_t)\}
\]
evolves along the trajectory. This is a geometric object that changes with
training, and it connects the observations without assuming a fixed group: early
training dominated by low-order statistical invariants, later training possibly
acquiring higher-level ones. It also carries an immediate falsifier --- if
$\operatorname{Stab}_t$ never grows on a corpus whose structure the model
demonstrably learns, then something stronger holds about strand generation than
any statement in this programme so far.

The prerequisite is a corpus on which the model learns structure at all. The
sweep of \S\ref{sec:stabilizer-evolution} failed that prerequisite, and the
suite --- symmetries exact by construction, histogram control, structural
ladder, $r$-stratification --- is in hand and unchanged for the rerun.

\subsection{What the programme now claims}
\label{sec:claims}

The disc is the slow state: optimiser memory, a diagonal Fisher metric,
accumulated statistics over $\Theta\times\mathcal{O}$. The strand is the fast
binary response field generated from the current batch, living on the vertices
of an $\ell^\infty$ ball, well defined only relative to the chart the disc
supplies. The dynamics respect the permutation gauge group of the
parameterisation, so the parameter index is a partition rather than a space.
The corpus acts on the strand through whatever invariants the current
representation has learned, and the symmetry structure is therefore a function
of training time rather than a property of the data.

Two things are established across corpora rather than on one: that the
batch~$\to$~strand map is high-dimensional (\S\ref{sec:corpus-generality}), and
that sparsity-based economy is refuted in consequence. The remaining
mathematical target is the smallest symmetry-respecting state space on which
disc and strand evolve --- and the remaining empirical question is unchanged
from \S\ref{sec:backward}: whether a dense substitution mechanism can generate
the strand more cheaply. Sparsity-based economy is refuted; substitution is
untested.
